% Include-ready complete bounds proofs through B19; generated from lvalue_bounds.md. Requires amsmath, amssymb, hyperref and a URL/path package.
\providecommand{\tightlist}{\setlength{\itemsep}{0pt}\setlength{\parskip}{0pt}}
\section{DEG\_L research: L-value bounds, fixed Weil projector, bounded witness theorem}\label{deg_l-research-l-value-bounds-fixed-weil-projector-bounded-witness-theorem}

Date: 21 November 2026. Actor: GPT-6 Astra Pro, subagent \texttt{/root/lvalue\_bounds}; effort setting not exposed. Run: DEG\_L\_20261121. This is a working research contribution to the root paper, not a claim to solve T6 or T7. Every assertion below is either explicitly conditional, an elementary derived finding with proof, or an identified cited background result. All degrees are geometric intersection degrees, without factorial normalization.

\subsection{Executive finding}\label{executive-finding}

The strongest positive reduction obtained here is \textbf{B8}: in a polarized family of abelian varieties of Weil type of dimension \(2n\), a uniformly bounded-degree effective \(n\)-cycle with \emph{any nonzero Weil component} supplies a bounded-division representative of every fixed flat rational Weil class. The multiplier and both signed degrees are \(O(D^2)\), with explicit constants depending only on the fixed family, polarization type, field action, and target class. No varying full CM field or varying Hecke projector is required. The proof gives a fixed polynomial correspondence in the generic imaginary-quadratic action, a comass estimate on invariant harmonic forms, and an integral two-dimensional Cramer calculation. This removes the denominator and isolation gates once bounded geometric witnesses exist; the existence of those witnesses in the open families remains an actual hypothesis.

The analytic findings are complementary. The genuine \(U(1)\to U(3)\) Rallis formula uses \(L_E(3/2,\chi)\), not \(L_E(1,\chi)\); its global factor has elementary conductor-independent upper/lower bounds for fixed field degree. It does not compute the degree of the unnormalized 032 image surface. Edge-of-strip bounds under the admitted Alpöge input and conditional Masullo input apply only after identifying the correct normalized \(L\)-function, special point, local vectors, measure, and CM extraction coefficients.

\subsection{Source findings and exact scope}\label{source-findings-and-exact-scope}

\textbf{S032.} OpenAI family 032, \emph{The rational Hodge conjecture for CM abelian varieties}, §4.4 Lemma 4.4 (printed pp.20--21) fixes finite Schwartz data and permits all sufficiently deep principal levels; §7 proof of Theorem 2.5, printed p.46, explicitly states that pulling to a deeper principal cover multiplies the mixed integral by the cover degree. The theorem then selects a nonzero term and rational Hodge maps by span and finite expansion; it supplies no estimate for coefficients of those maps. These are statements used and rechecked in this run, not independent reproofs of 032.

\textbf{S-Rallis1.} Jaeho Haan, \emph{The Bessel Period of U(3) and U(2) involving a non-tempered representation}, arXiv:1403.5061v7, §3.3, Theorem 3.4, printed p.8, gives, in that paper's Tamagawa and splitting conventions, \[
\langle\theta(\bar f_1,\varphi_1),\theta(\bar f_2,\varphi_2)\rangle
 = \frac{L_E(3/2,\operatorname{BC}(\tau)\otimes\gamma^3)}{L_F(3,\epsilon_{E/F})}
 \prod_v Z_v^{\#}(1,f_{1,v},f_{2,v},\Phi_{1,v}).
\] Here \(\tau\) is an automorphic character of \(U(1)\) and \(\gamma\) is the chosen unitary splitting character; §3.1 (3.1)--(3.2) defines the normalization of local zeta integrals. §3.6.1, printed p.9, explicitly records absolute convergence and nonvanishing of the numerator at \(3/2\). This is a formula for integrated theta lifts with source character, and cannot be substituted for an arbitrary theta-kernel evaluation or wedge in 032 without comparison.

URL: https://arxiv.org/pdf/1403.5061 .

\textbf{S-Rallis2.} Gan--Qiu--Takeda, \emph{The Regularized Siegel--Weil Formula (The Second Term Identity) and the Rallis Inner Product Formula}, arXiv:1207.4709v3, Theorem 11.4, printed pp.54--55, under its explicit range and cuspidality/first-occurrence hypotheses, gives the unitary point \((m-n+1)/2\), with twisting character and normalized doubling integrals. Its range \(n<m\leq2n\) excludes \(U(1)\to U(3)\). The \(U(2)\to U(3)\) point is \(1\), but the standard degree-two \(L\)-function is not thereby a Hecke-character \(L\)-function. A source spectral decomposition or endoscopic factorization must be supplied. The theorem is used as background and its range checked, not reproved.

URL: https://arxiv.org/pdf/1207.4709 .

\textbf{S-Alpöge.} The user's Dirichlet GRH input is admitted. Exact supplied Lean declarations: \texttt{ALPOGE\_Main.lean}, \path|LFunction_eq_zero_re_of_isPrimitive : PrimitiveDirichletLZerosOnLine|, and \path|LFunction_ne_zero_of_half_lt_re : DirichletLNonvanishingRightOfHalf|; definitions in \texttt{ALPOGE\_Statements.lean}. The supplied \emph{Remark}, step 20, printed p.4, announces GRH; steps 9, 12, 18 concern lift norms, Rallis, and indices respectively. The full imported archive is not supplied, and this run did not compile it. No claimed new divisibility theorem is inferred from those steps or from the mapping file.

\textbf{S-Masullo.} The full draft is not supplied. The prompt attributes Hecke GRH to §1.2. The supplied \texttt{MASULLO\_STATEMENTS\_INDEX.md} is an index of openings, not theorem statements or proofs. It records Chapter 12, Theorem 12.5.1.1, PDF p.519, and Chapter 13, Theorem 13.3.1.1, PDF p.522, as trace formulas, and records an unfinished-looking opening labelled 13.3.1.2 (at the text-layer level). None is silently promoted to a precise theorem of Hecke GRH in this report. All use of Hecke GRH is conditional on the claim specified by the user, not certified from the index.

\subsection{B1. Elementary uniform bounds in the absolutely convergent half-plane}\label{b1.-elementary-uniform-bounds-in-the-absolutely-convergent-half-plane}

\textbf{Statement.} Let \(K\) have degree \(d\), and let \(\chi\) be a unitary Hecke character. For real \(\sigma>1\), using the finite Euler product with ramified factors equal to one, \[
 \left(\frac{\zeta(2\sigma)}{\zeta(\sigma)}\right)^d
 \leq |L_K(\sigma,\chi)|\leq \zeta(\sigma)^d.
\] More sharply, \(\zeta_K(2\sigma)/\zeta_K(\sigma)\leq|L_K(\sigma,\chi)|\leq\zeta_K(\sigma)\). No RH or GRH is needed.

\textbf{Proof.} For every unramified prime ideal \(\mathfrak p\), unitarity gives \(|\chi(\mathfrak p)|=1\). Thus \[
 (1+N\mathfrak p^{-\sigma})^{-1}\leq
 |1-\chi(\mathfrak p)N\mathfrak p^{-\sigma}|^{-1}
 \leq(1-N\mathfrak p^{-\sigma})^{-1}.
\] The same inequalities hold with the middle term replaced by one at ramified primes. Absolute convergence permits multiplication. The left product is \(\zeta_K(2\sigma)/\zeta_K(\sigma)\). Above each rational prime there are at most \(d\) prime ideals, each with norm at least that prime. Comparing each factor gives the stated bounds by rational zeta functions. All factors are positive and the products converge. QED.

\textbf{Application to the finite Euler factors.} In S-Rallis1 with \([F:\mathbb Q]=d\) and \([E:\mathbb Q]=2d\), the absolute value of the ratio of the finite Euler products lies between \[
 \frac{\zeta(3)^d}{\zeta(3/2)^{2d}}
 \quad\hbox{and}\quad
 \zeta(3/2)^{2d}\left(\frac{\zeta(3)}{\zeta(6)}\right)^d.
\] This follows by applying B1 separately to numerator and denominator. It controls no local \(Z_v^\#\) and no change of measure, input scale, or cycle normalization. When the \(L\)-symbols include their archimedean local factors, those factors must also be included; the actual cohomological situation is evaluated next.

\subsubsection{B1a. Actual infinity type and the location of the rank-one value}\label{b1a.-actual-infinity-type-and-the-location-of-the-rank-one-value}

For an actual \(H^{1,0}\) oscillator summand, choose the target splitting character \(\gamma=\mu\) from the supplied van Hoften--Sempliner §3.1.3, printed p.9. It has infinity type \(\gamma_v(z)=(z/|z|)^{\varepsilon_v}\) with \(\varepsilon_v\in\{1,-1\}\). Retwist the source as in T6.11 to Haan's compatible pair \((\gamma^3,\gamma)\). Then the source character \(\tau\) is trivial at every archimedean place. Consequently \[
 \chi=\operatorname{BC}(\tau)\gamma^3,
 \qquad \chi_v(z)=(z/|z|)^{3\varepsilon_v},
 \qquad L_v(s,\chi_v)=\Gamma_{\mathbb C}(s+3/2).
\tag{B1a.1}
\] Here \(\Gamma_{\mathbb C}(s)=2(2\pi)^{-s}\Gamma(s)\) and \(\Gamma_{\mathbb R}(s)=\pi^{-s/2}\Gamma(s/2)\).

\textbf{Proof.} The center of \(U(3)\) acts trivially on its symmetric domain and on its cotangent representation; at the other real places the differential-form coefficient has the trivial compact representation. Thus every target archimedean center acts trivially on the actual cohomological vector and on its generated representation. Haan, §2.1, Remark 2.1, printed p.4, states that for a one-dimensional member the images of that member and of the center of the other member coincide in the metaplectic group under his power-compatible splittings. Equivalently their scalar actions in a nonzero local theta quotient match, possibly after inversion according to the theta convention. It follows that the source \(U(1)\) character is trivial. Its base change \(z\mapsto\tau_v(z/\bar z)\) is therefore trivial as well. The displayed infinity type follows. The standard complex local factor can also be read directly from the Gaussian Mellin integral: after canceling the angular exponent by the appropriate polynomial, its radial integral is a nonzero fixed multiple of \(\int_0^\infty r^{2s+|3\varepsilon_v|-1}e^{-2\pi r^2}\,dr\), giving \(\Gamma_{\mathbb C}(s+3/2)\). QED.

Since \(\epsilon_{E/F,v}=\operatorname{sgn}\) at every real place of \(F\), its local factor at \(3\) is \(\Gamma_{\mathbb R}(4)\). Thus the all-place product of archimedean ratios in the rank-one formula is exactly \[
 \left(\frac{\Gamma_{\mathbb C}(3)}{\Gamma_{\mathbb R}(4)}\right)^d
 =(2\pi)^{-d}.
\tag{B1a.2}
\] With these explicit local \(L\) conventions, multiply both finite-product bounds above by \((2\pi)^{-d}\) when using all-place \(L\)-functions. This is a fixed factor, independent of the finite conductor. A separately completed functional equation may put a further \(q^{s/2}\) into a symbol \(\Lambda(s)\); that completion is not the all-place local product appearing in the Rallis normalization and must not be inserted silently.

The unitary functional-equation center is \(1/2\), while the Rallis point is \(3/2\), in the absolutely convergent half-plane. If \(\gamma=\psi\|\cdot\|_E^{1/2}\) with the algebraic character \(\psi\) supplied in the proof of van Hoften--Sempliner Lemma 3.1.5, pp.9--10, put \(\Psi=\operatorname{BC}(\tau)\psi^3\). Then \[
 L(s,\chi)=L(s+3/2,\Psi).
\] The arithmetic normalization of weight three has center \(2\), and the Rallis value becomes \(L(3,\Psi)\). Its gamma factors and those of the reflected conjugate are \(\Gamma_{\mathbb C}(s)\) and \(\Gamma_{\mathbb C}(4-s)\), respectively, so the critical integers are \(1,2,3\). Thus \(3\) is the rightmost critical integer, not the central one. This arithmetic description is a normalization statement; the unconditional bound uses the unitary finite Euler product at \(3/2\).

\subsection{B2. Exact edge bounds and what each GRH input buys}\label{b2.-exact-edge-bounds-and-what-each-grh-input-buys}

Lamzouri--Li--Soundararajan, \emph{Conditional bounds for the least quadratic non-residue and related problems}, arXiv:1309.3595v3, Theorem 1.5 (printed pp.3--4), gives for primitive Dirichlet characters of conductor \(q\geq10^{10}\) under GRH, with \[
 A(q)=\log\log q-\log2+\tfrac12+1/\log\log q,
\] \[
 |L(1,\chi)|\leq 2e^\gamma A(q),\qquad
 |L(1,\chi)|^{-1}\leq\frac{12e^\gamma}{\pi^2}
 \left(A(q)+\frac{14\log\log q}{\log q}\right).
\] The constant in the second inequality is \(12e^\gamma/\pi^2\), not \(12e^{2\gamma}/\pi^2\). An older hosted text extraction returned the latter; the latest arXiv version and Lumley's Theorem 1.1 provide the correct former value. Small conductors require their own finite verification or a separate bound. Under the user's admitted Alpöge theorem, this consequence no longer has a GRH condition. URL: https://arxiv.org/pdf/1309.3595.

For primitive nontrivial unitary Hecke characters whose completed \(L\)-function is entire, write the completed function in degree \(d=[K:\mathbb Q]\) over \(\mathbb Q\), with conductor \(q(\chi)=|D_K|N\mathfrak f_\chi\) and archimedean parameters \(\kappa_j\). Allysa Lumley, \emph{Explicit bounds for L-functions on the edge of the critical strip}, arXiv:1804.09850v1, §1.1 and Theorems 2.1--2.2 (printed pp.3--4), applies to \[
 C(\chi)=\frac{q(\chi)}{\pi^d}\prod_{j=1}^d\left|\frac{1+\kappa_j}{2}\right|.
\] For fixed degree, its bounds have orders \((\log\log C)^{d}\) and \((\log\log C)^{-d}\); its explicit formulas apply when \(\log C\geq23d\). Unit modulus of unramified Hecke parameters supplies the required Ramanujan hypothesis; the remaining GRH input is conditional on Masullo for general \(K\). Pole cases (including the trivial character) are excluded. Field degree and archimedean type cannot be hidden while allowing them to vary. URL: https://arxiv.org/pdf/1804.09850.

These are known results applied, not new theorems. In particular, this report does not assert an edge bound for a nonabelian \(U(2)\) standard \(L\)-function from Hecke GRH. Nor does GRH give a positive lower bound at the center: the functional equation can force a central zero when a self-dual root number is \(-1\).

\subsection{B3. Level, vector scale, and arithmetic conductor are distinct}\label{b3.-level-vector-scale-and-arithmetic-conductor-are-distinct}

\textbf{Statement.} If \(\pi:Y\to X\) is a finite étale cover of degree \(d\) and \(\eta\) is a top-degree form on \(X\), then \(\int_Y\pi^*\eta=d\int_X\eta\). Consequently a nonzero pulled-back theta form has unnormalized Petersson norm squared multiplied by \(d\). If its underlying automorphic representation, infinity type, and finite vector are fixed, the representation's arithmetic and analytic conductors stay fixed.

\textbf{Proof.} Integration over evenly covered coordinate sets and a partition of unity gives the integral identity. Apply it to the positive top form defining the norm. Conductors are invariants of the fixed representation and archimedean parameters, not of which smaller compact open subgroup happens to fix its vector. QED.

This is also explicitly recorded for the mixed period in 032 §7 p.46. It is an exact level-growth statement for a chosen presentation, not a lower bound for the intrinsic minimum complexity \(c(s)\) in Route 2. Similarly scaling a Schwartz input by \(u\) scales a theta lift by \(u\) and its norm squared by \(|u|^2\), while leaving its representation's \(L\)-function unchanged. The local integrals and/or input pairing carry that factor.

\subsection{B4. Polarized isogenies: rational degree can stay fixed while denominators grow}\label{b4.-polarized-isogenies-rational-degree-can-stay-fixed-while-denominators-grow}

\textbf{Statement.} Let \(\psi:A\to B\) be an isogeny of \(g\)-dimensional polarized abelian varieties. Suppose \(\psi^*c_1(L_B)=c\,c_1(L_A)\) and \(L_A,L_B\) have the same polarization degree. For a codimension-\(p\) cycle \(Z\), \[
 \deg\psi=c^g,\qquad
 \deg_{L_A}(\psi^*Z)=c^p\deg_{L_B}Z.
\]

\textbf{Proof.} Integrating the \(g\)th power of the polarization identity gives the first equality. By projection formula, \[
 \int_A\psi^*[Z]c_1(L_A)^{g-p}
 =c^{-(g-p)}\deg\psi\int_B[Z]c_1(L_B)^{g-p},
\] which is the second. QED.

The rational cycle \(c^{-p}\psi^*Z\) has unchanged degree but its displayed coefficient has denominator \(c^p\) when \(c\) is an integer. For multiplication by \(m\) on a fixed variety, \(c=m^2\). This is not a proof that all representatives have large denominators: it is a precise diagnosis of the natural pullback representative.

\subsection{B5. A rigorous norm-to-degree bridge for surface maps}\label{b5.-a-rigorous-norm-to-degree-bridge-for-surface-maps}

\textbf{Statement.} Let \(Y\) be a smooth projective surface, \(\omega\) a Kähler class, \(V=\int_Y\omega^2\), and \(f:Y\to A\) a morphism with two-dimensional image \(Z\) and generic degree \(\delta\). Choose a translation-invariant representative \(\omega_A\in c_1(L)\) of a polarization, and invariant holomorphic forms \(\alpha_1,\ldots,\alpha_g\) such that \(\omega_A=i\sum_j\alpha_j\wedge\bar\alpha_j\). Put \[
 E_f=\int_Y f^*\omega_A\wedge\omega
 =\sum_j\|f^*\alpha_j\|_\omega^2,\qquad
 \|u\|_\omega^2=i\int_Yu\wedge\bar u\wedge\omega.
\] Then \[
 \deg_L Z=\frac{\int_Y(f^*\omega_A)^2}{\delta}
 \leq\frac{E_f^2}{\delta V}.
\]

\textbf{Proof.} Proper pushforward gives \(f_*[Y]=\delta[Z]\) and projection formula gives the equality. With \(\beta=f^*\omega_A\), decompose \(\beta=(E_f/V)\omega+\beta_0\), where \(\int\beta_0\omega=0\). The Hodge index theorem on the surface gives \(\int\beta_0^2\leq0\), so \(\int\beta^2\leq E_f^2/V\). The displayed energy identity follows by linearity from the specified polarization representative. QED.

This is a direct application of the classical surface Hodge index theorem in a form adapted to the present problem, not a claim to reprove that theorem. Under a cover of degree \(d\), \(E_f\) and \(V\) both scale by \(d\). Thus the bound \(E_f^2/V\) for the degree of the pushforward class scales by \(d\). For the reduced image, the generic mapping degree \(\delta\) also scales by \(d\), and the displayed bound remains unchanged; this matches B3 exactly.

For the sufficiently deep principal ball-quotient surfaces used in 032, the actual pullbacks \(f^*\alpha_j\) do have finite expansions \(\sum_a b_{ja}\theta_a\) in genuine integrated theta lifts from characters of one-dimensional unitary groups. The existence of these expansions is supplied by Bergeron--Millson--Mœglin, \emph{The Hodge conjecture and arithmetic quotients of complex balls}, Acta Math.216 (2016), Theorem 9.1, pp.73--75, with the range in §9.1, p.72: for \(m=3\) and \((a,b)=(0,1)\), its condition is \(4<6\). The supplied van Hoften--Sempliner paper makes the all-\(H^{1,0}\) spanning conclusion explicit in the proof of Proposition 4.1.8, pp.15--16; Proposition 3.1.7, p.11, identifies the oscillator constituents. The published BMM proof and the supplied Proposition 4.1.8 proof were read. This is imported classification/spanning, not a new proof of those theorems.

For completeness, only finitely many direct-limit theta vectors are needed for each fixed form. Put them at a common finer level and average over the desired compact open subgroup. The original form is invariant, and equivariance moves the finite average to the Schwartz inputs, so the expansion exists at that level. Equality in \(H^{1,0}\) is equality of holomorphic forms by the injectivity into Betti cohomology for a compact Kähler variety. If necessary one first deepens the principal level so the arithmetic \(U\) and \(SU\) groups agree: their determinant values are norm-one algebraic units, hence roots of unity, and sufficiently deep congruence kills the finite set of nontrivial roots. T6.9 gives the full argument and its scope.

With norm squares \(Q_a\) computed by a valid rank-one Rallis formula, and with the scalar/geometric metric and measure conversions included in \(Q_a\), the triangle inequality gives the concrete bound \[
 \deg_L Z\leq\frac{1}{\delta V}
 \left\{\sum_j\left(\sum_a|b_{ja}|\sqrt{Q_a}\right)^2\right\}^2.
\] The existence gate is therefore removed. What remains is a quantitative bound on the coefficients \(b_{ja}\) relative to the integral Hom lattice, together with CM periods, test-vector factors and metric/measure normalization. Orthonormalizing the theta basis does not control those integral coefficients. B10 below reduces them to an integral energy-lattice covolume estimate, and B11 shows exactly why that estimate cannot be inferred from finite generation or orthonormalization alone.

\subsection{B6. Fixed generic-field projector onto Weil classes}\label{b6.-fixed-generic-field-projector-onto-weil-classes}

\textbf{Setup.} Let \(A\) be a polarized abelian variety of dimension \(2n\), carrying an action of an order \(\mathcal O\subset K\) in an imaginary quadratic field. Assume the Rosati involution restricts to conjugation on \(K\) and the two eigenspaces in \(H^{1,0}\) both have dimension \(n\) (Weil type). Let \[
 W=\bigwedge_K^{2n}H^1(A,\mathbb Q)\subset H^{2n}(A,\mathbb Q).
\] Here the inclusion means the rational rank-two subspace whose complexification is the sum of the two extreme \(K\)-eigenlines. Choose a nonzero \(a\in\mathcal O\) such that \(a/\bar a\) is not a root of unity, and let \(T=[a]^*\) on cohomology. Such an \(a\) exists: only finitely many arguments arise from roots of unity in \(K\), whereas the arguments of nonreal elements \(M+\tau\) for a fixed nonreal \(\tau\in\mathcal O\) vary with integers \(M\) and tend to zero.

\textbf{Statement.} There is a polynomial \(P(X)\in\mathbb Z[X]\) and a positive integer \(R\), depending only on \((K,n,a)\), such that \(e=P(T)/R\) is the projector onto \(W\) on \(H^{2n}\). It is induced by a fixed rational algebraic correspondence and is self-adjoint for the Poincaré pairing on that degree. Its denominator is independent of the CM point or Hecke level. One possible choice is \[
 \begin{aligned}
 \lambda_r&=a^r\bar a^{2n-r}\qquad(0\leq r\leq2n),\\
 f_W(X)&=(X-a^{2n})(X-\bar a^{2n}),\\
 f_R(X)&=\prod_{r=1}^{2n-1}(X-\lambda_r),\\
 R&=\operatorname{Res}(f_W,f_R)>0,
 \end{aligned}
\] with \(P=vf_R\) for integral polynomials satisfying \(u f_W+v f_R=R\).

\textbf{Proof.} Decomposing \(H^1_\mathbb C=V_a\oplus V_{\bar a}\), each summand has dimension \(2n\), and \[
 H^{2n}_\mathbb C=\bigoplus_{r=0}^{2n}\bigwedge^rV_a\otimes\bigwedge^{2n-r}V_{\bar a}.
\] The corresponding eigenvalues are \(\lambda_r\). They are distinct since \(a/\bar a\) has infinite order. Each polynomial displayed is fixed by conjugation and has algebraic-integer coefficients, hence lies in \(\mathbb Z[X]\); the two factors are coprime. The Sylvester matrix for these monic polynomials has determinant their resultant. Multiplying that matrix by its integral adjugate and extracting the column corresponding to the constant polynomial gives integral \(u,v\) with \(u f_W+v f_R=R\) (a simultaneous sign change adjusts the determinant convention). Thus \(P/R\) takes value one at the two extreme eigenvalues and zero at every other eigenvalue. Since \(T\) is diagonalizable, it is the asserted projector.

Every power \(T^j\) is induced by pullback along the algebraic isogeny \([a^j]\), so the polynomial defines an algebraic correspondence. Its cohomological projector property is what is asserted: \textbf{no Chow-group identity \(e^2=e\) is claimed} merely from this calculation. For a full homological-motive statement, compose it with the algebraic degree projector \(\pi_{2n}=\prod_{j\ne2n}([2]^*-2^j)/(2^{2n}-2^j)\), where \(0\leq j\leq4n\); this acts as the identity on \(H^{2n}\) and as zero in all other degrees. Thus it defines an idempotent on \(h^{2n}\) in homological motives, and on numerical cycle classes, without claiming a Chow idempotence relation.

To prove self-adjointness on this degree, the Poincaré pairing of two \(T\)-eigenvectors can be nonzero only if \(\lambda_r\lambda_s=\deg[a]=N(a)^{2n}\): this follows by applying \([a]^*\) to their product in top degree. Infinite order of \(a/\bar a\) makes this condition exactly \(r+s=2n\). Consequently the two extreme eigenlines pair only with each other and are orthogonal to all the intervening eigenspaces. Projection onto their sum is Poincaré-self-adjoint. Finally, \(R=f_R(a^{2n})\overline{f_R(a^{2n})}>0\). QED.

A useful explicit expression is \[
 R=N(a)^{n(2n-1)}\prod_{k=1}^{2n-1}|a^k-\bar a^k|^2.
\] It follows by factoring \(f_R(a^{2n})=a^{n(2n-1)}\prod_{k=1}^{2n-1}(a^k-\bar a^k)\) and taking the field norm.

If \(P=\sum_jb_jX^j\), put \[
 B_P=\sum_j|b_j|N(a)^{jn}.
\] For effective codimension-\(n\) cycles \(Z\) of degree \(D\), \(P([a]^*)Z\) is an integral signed cycle whose positive and negative parts both have degree at most \(B_PD\). Indeed \([a]^*c_1(L)=N(a)c_1(L)\) by the Rosati condition, so B4 gives degree multiplier \(N(a)^n\) for pullback on codimension \(n\); use it term by term. This works fibrewise in any family with the fixed action and polarization data.

\subsection{B7. Uniform Weil-component bound from effective degree}\label{b7.-uniform-weil-component-bound-from-effective-degree}

\textbf{Statement.} In B6 put \(\omega=c_1(L)\), \(V_A=\int_A\omega^{2n}/(2n)!=\chi(L)\), and \[
 q(u,v)=(-1)^n\int_Au\wedge v\quad(u,v\in W_\mathbb R).
\] Then \(q\) is positive definite. For every effective \(n\)-cycle \(Z\) of degree \(D=\int_Z\omega^n\), \[
 \|e[Z]\|_q\leq\frac{D}{n!\sqrt{V_A}}.
\] The constant depends only on polarization type and \(n\), hence is uniform on the whole Weil family, without a compactness restriction on its base.

\textbf{Proof.} The Rosati identity makes \(\omega\) a mixed-eigenvalue class in \(V_a\otimes V_{\bar a}\): its \([a]^*\) eigenvalue is \(N(a)\) on degree two, different from \(a^2\) and \(\bar a^2\). Wedge product by \(\omega\) kills each extreme degree-\(2n\) eigenline because it would require \(2n+1\) vectors in one \(2n\)-dimensional eigenspace. Thus \(W\) is primitive. Weil type makes both lines of Hodge type \((n,n)\). The Hodge--Riemann bilinear relations give positivity of \(q\); the Hodge star on these primitive middle-degree real classes is \(*=(-1)^n\).

Give \(A\) its flat translation-invariant metric associated with \(\omega\). Every harmonic form is translation invariant, so for its harmonic representative \(\eta\in W_\mathbb R\), the pointwise Euclidean norm is constant and equals \(\|\eta\|_q/\sqrt{V_A}\). Its comass is at most its Euclidean norm. The complex analytic cycle \(Z\) has volume \(D/n!\) by Wirtinger's identity, with multiplicities; its singular locus has measure zero. Therefore \[
 \left|\int_Z\eta\right|\leq\frac{D}{n!\sqrt{V_A}}\|\eta\|_q.
\] By B6, \(e\) is self-adjoint for Poincaré pairing, so for \(w=e[Z]\), \[
 \|w\|_q^2=|q(w,w)|=\left|\int_A[Z]\wedge w\right|=\left|\int_Zw\right|.
\] If \(w=0\) the result holds; otherwise apply the preceding inequality with \(\eta=w\) and divide by \(\|w\|_q\). QED.

\subsection{B8. Effective witness implies bounded division, with quadratic bounds}\label{b8.-effective-witness-implies-bounded-division-with-quadratic-bounds}

\textbf{Statement.} Keep a connected polarized Weil-type family as in B6, with a globally defined \(\mathcal O\)-action (after finite base change if necessary), fixed \(n,K,a,L\), and fixed rational flat target \(t\in W\). Put \[
 \Lambda=W\cap H^{2n}(A,\mathbb Z),\qquad \delta_\Lambda=\det(q(e_i,e_j))>0
\] for any integral basis of this rank-two saturated lattice. Choose \(b\geq1\) with \(\tau=bt\in\Lambda\), and write \(h=N(a)^n\). The constants \(V_A\), \(\delta_\Lambda\), and \(\|\tau\|_q\) are locally constant along the family and so fixed on a connected base.

Suppose a fibre contains an effective \(n\)-cycle \(Z\) of degree \(D\) with \(e[Z]\ne0\). Then there are an integer \(N\geq1\) and effective \(Z^+,Z^-\) with \[
 Nt=[Z^+]-[Z^-],
\] where, with \(M=RD/(n!\sqrt{V_A})\), \[
 N\leq\frac{b h M^2}{\sqrt{\delta_\Lambda}},\qquad
 \deg Z^\pm\leq\frac{2h\|\tau\|_q M B_P D}{\sqrt{\delta_\Lambda}}.
\] Thus both are \(O(D^2)\), uniformly in the fibre. In particular a uniform bound for degrees of nonzero Weil witnesses on a Zariski-dense set gives bounded division on that set. The proper parameter-space argument of Route 2 then applies, once supplied separately with its stated hypotheses.

\textbf{Proof.} Let \(V_Z=P([a]^*)Z\) as an integral signed cycle and \(v=[V_Z]=R e[Z]\). This is a nonzero element of \(\Lambda\), with \(\|v\|_q\leq M\) by B7. The characteristic polynomial of \(T\) on \(W\) is \(f_W\). Since \(a/\bar a\) is not a root of unity, \(a^{2n}\) is nonreal, so \(f_W\) is irreducible over \(\mathbb Q\). Hence \(v,Tv\) are a rational basis of \(W\): otherwise \(v\) would be a rational eigenvector of \(T\) with a rational eigenvalue.

Use any oriented integral basis of \(\Lambda\). Its determinant \[
 \Delta=\det(v,Tv)\in\mathbb Z\setminus\{0\}
\] therefore satisfies \(\Delta\tau=m v+\ell Tv\), where the integral Cramer numerators are \(m=\det(\tau,Tv)\) and \(\ell=\det(v,\tau)\). On \(W\), \(T\) is a similitude of \(q\) of multiplier \(\deg[a]=N(a)^{2n}=h^2\), so \(\|Tv\|_q=h\|v\|_q\). The area of a parallelogram in this metric is \(\sqrt{\delta_\Lambda}\) times its lattice determinant. The elementary area bound consequently gives \[
 |\Delta|\leq hM^2/\sqrt{\delta_\Lambda},\quad
 |m|\leq\|\tau\|_q hM/\sqrt{\delta_\Lambda},\quad
 |\ell|\leq\|\tau\|_q M/\sqrt{\delta_\Lambda}.
\] Let \(N=b|\Delta|\), and let \(C=\operatorname{sgn}(\Delta)(mV_Z+\ell[a]^*V_Z)\). Then \([C]=Nt\). By the final part of B6 the positive and negative degrees of \(V_Z\) are at most \(B_PD\), and those of \([a]^*V_Z\) are at most \(hB_PD\). Splitting the signed cycle \(C\) into positive and negative parts thus bounds each by \[
 (|m|+h|\ell|)B_PD\leq 2h\|\tau\|_q M B_PD/\sqrt{\delta_\Lambda}.
\] This proves the claim. The intersection lattice and its pairing vary as an integral local system, so its discriminant and the norm of a fixed flat target are constant; \(V_A\) is fixed by polarization type. No choice of a global basis is needed to make the constants uniform. QED.

\textbf{Stronger exact index identity.} If \(\mu=a^{2n}\) under the fixed complex embedding, then \[
 |\Delta|=\frac{|\operatorname{Im}\mu|}{\sqrt{\delta_\Lambda}}\,\|v\|_q^2.
\] Indeed \(T\) is multiplication by the nonreal scalar \(\mu\) on the real plane \(W\), its characteristic polynomial has trace \(2\operatorname{Re}\mu\), and it is a positive-orientation \(q\)-similitude. In a \(q\)-orthonormal basis it is the rotation-dilation matrix for \(\mu\) or \(\bar\mu\); the area of \((v,Tv)\) is \(|\operatorname{Im}\mu|\|v\|^2\). Dividing by the lattice covolume proves the formula. This makes the last arithmetic index precise, without claiming it is automatically an automorphic \(L\)-value.

\textbf{The quadratic size is sharp for extraction from one lattice vector.} In the positive lattice \(\Lambda=\mathbb Z[i]\), let \(v=m+i\) and target \(\tau=1\). If \(N\tau\) belongs to the principal ideal \(v\mathbb Z[i]\), then \[
 \frac{N}{m+i}=\frac{Nm}{m^2+1}-i\frac{N}{m^2+1}\in\mathbb Z[i],
\] so \(m^2+1\mid N\). Since \(\|v\|^2=m^2+1\), extraction using \(v\) and any fixed integral multiplication operator can require an integer of quadratic size in \(\|v\|\). This is an arithmetic sharpness statement for that extraction method. It does not establish growth of intrinsic cycle complexity, since other geometric cycles might represent the target more efficiently.

\textbf{What remains.} The theorem does not create \(Z\) in the non-CM fibres, and 032 gives neither a uniform degree bound for such \(Z\) on a dense orbit nor a deformation of one actual full cycle class. It proves that a separate unbounded full-CM-field projector is unnecessary in the Weil-family route. Once geometric witness degrees are bounded, the two-dimensional integral extraction is automatically bounded.

\subsection{B9. A precise failure of archimedean norm control for arbitrary projectors}\label{b9.-a-precise-failure-of-archimedean-norm-control-for-arbitrary-projectors}

On the lattice \(\mathbb Z^2\) with Euclidean pairing, the orthogonal projector to \(\mathbb Q(m,1)\) is \[
 E_m=\frac1{m^2+1}\begin{pmatrix}m^2&m\\m&1\end{pmatrix}.
\] It has operator norm one, while its exact common denominator is \(m^2+1\), because its lower-right entry is \(1/(m^2+1)\). The calculation follows by the rank-one projector formula \(vv^t/(v^tv)\). This example already lives in \(\operatorname{End}(\mathbf1\oplus\mathbf1)\) in homological or numerical motives: all matrices are algebraic correspondences between two copies of a point. It refutes any abstract inference from bounded archimedean projector norm to bounded integral denominator. It does not refute B6, which fixes the generic-field projector, or the existence of other bounded representatives in a geometric family.

Similarly \(x_m=(1,1/m)\) has norm in \([1,\sqrt2]\) while its pairing with \((0,1)\) is \(1/m\). Lower bounds for two individual norms therefore give no lower bound for the mixed period needed in 032. The positive continuous comparison used by 032 is a separate input, whose quantitative approximation costs must be controlled.

\subsection{B10. An integral Hom covolume bound replaces individual coefficient bounds}\label{b10.-an-integral-hom-covolume-bound-replaces-individual-coefficient-bounds}

\textbf{Setup.} Let \(Y\) be a smooth projective variety of dimension \(d\ge1\), with an integral ample class \(H\), and fix a base point for its Albanese map \(\alpha:Y\to B=\operatorname{Alb}(Y)\). Let \((A_i,L_i)\) be polarized abelian varieties, for \(1\le i\le k\). Put \[
 \Lambda_i=\operatorname{Hom}(B,A_i),\qquad
 E_i(u)=\int_Y (u\alpha)^*c_1(L_i)H^{d-1},\qquad
 V=\int_YH^d.
\] Here maps are homomorphisms; translations do not affect the forms, energies or degrees. Assume \(r_i=\operatorname{rank}\Lambda_i>0\). Let \(q_i\) be the real quadratic metric with square norm \(E_i\), let \(C_i\) be the Euclidean covolume of \(\Lambda_i\) for \(q_i\), and let \(\mu_i>0\) be any known lower bound for \(E_i(u)\) for nonzero integral \(u\). One may always take \(\mu_i=1\). Set \[
 v_r=\frac{\pi^{r/2}}{\Gamma(1+r/2)},\qquad
 R_i=\frac{2^{r_i}}{v_{r_i}}\frac{C_i}{\mu_i^{(r_i-1)/2}}.
\]

\textbf{Statement.} The lattice \(\Lambda_i\) contains \(r_i\) independent homomorphisms \(u_{ij}\), spanning \(\Lambda_i\otimes\mathbb Q\), such that \(E_i(u_{ij})\le R_i^2\). Suppose a specified real-multilinear period \(P\) on \(\prod_i(\Lambda_i\otimes\mathbb R)\) is nonzero. Some tuple of these homomorphisms has \(P(u_{1j_1},\ldots,u_{kj_k})\ne0\). If the period is the integral of a wedge of \(2d\) pulled-back one-forms, its nonvanishing implies that the product map \[
 F=(u_{1j_1}\alpha,\ldots,u_{kj_k}\alpha):Y\longrightarrow\prod_i A_i
\] has a \(d\)-dimensional image \(Z\). For the product polarization \(L=\bigotimes_i\operatorname{pr}_i^*L_i\), and the actual generic degree \(\delta_F\) of \(F\) onto \(Z\), \[
 \deg_L Z\le\frac{(\sum_iR_i^2)^d}{\delta_F V^{d-1}}.
\tag{B10.1}
\] For 032, \(d=2\) and \(k=4\). More generally, if a target polarization \(L'\) satisfies \(B c_1(L)-c_1(L')\) nef for a known \(B>0\), the right side is multiplied by \(B^d\).

\textbf{Proof.} Homomorphisms inject into homomorphisms of their integral first homology lattices, so \(\Lambda_i\) is a finitely generated free abelian group. Pullback on invariant one-forms depends linearly on the homomorphism. Writing a polarization as a positive sum of \(i\eta\wedge\bar\eta\) shows that \(E_i\) is a positive semidefinite real quadratic form. There is no nonzero integral vector in its kernel: if all pulled-back one-forms vanish on \(Y\), the map \(u\alpha\) has zero derivative and is constant; because \(\alpha(Y)\) generates \(B\) and both maps preserve the chosen origins, \(u=0\). The intersection expression for \(E_i\) is integral, so polarization gives a rational Gram matrix in every integral basis. A nonzero real kernel of a rational matrix contains a nonzero rational vector, and clearing its denominator would give a nonzero integral kernel vector. Thus the real kernel is zero and \(q_i\) is positive definite. Its positive values on nonzero lattice vectors are at least one.

Let \(\lambda_1\le\cdots\le\lambda_{r_i}\) be the successive minima of the Euclidean unit ball for this metric. Minkowski's second theorem gives \[
 v_{r_i}\prod_j\lambda_j\le 2^{r_i}C_i.
\] Since every \(\lambda_j\ge\sqrt{\mu_i}\), one obtains \(\lambda_{r_i}\le R_i\). The minima are attained, giving the required independent maps. They need not be an integral basis, and none is asserted: a rational spanning set is enough. If \(P\) vanished on every tuple from these spanning sets, multilinearity would make it vanish identically, a contradiction. A map with image dimension smaller than \(d\) pulls every smooth real \(2d\)-form back to zero, so the nonzero wedge period gives the image-dimension assertion.

For the degree bound, write \(\beta=F^*c_1(L)\), a nef class. Its energy is \(E=\beta H^{d-1}=\sum_iE_i(u_{ij_i})\le\sum_iR_i^2\). The mixed-intersection inequality \[
 \beta^d\le\frac{(\beta H^{d-1})^d}{(H^d)^{d-1}}
\tag{B10.2}
\] follows directly from surface Hodge index: for rational ample \(\beta,H\), take smooth complete-intersection surfaces of suitable very ample multiples of \(j-1\) copies of \(\beta\) and \(d-j-1\) copies of \(H\), for \(1\le j\le d-1\). Hodge index on each surface gives \(a_j^2\ge a_{j-1}a_{j+1}\) for \(a_j=\beta^jH^{d-j}\). These inequalities imply \(a_d\le a_1^d/a_0^{d-1}\). Approximating nef \(\beta\) by \(\beta+\varepsilon H\) proves the required nef case by continuity. For \(d=1\) it is equality, and for \(d=2\) it is B5. Projection formula gives \(\beta^d=\delta_F\deg_L Z\), proving (B10.1). Finally nefness of \(Bc_1(L)-c_1(L')\) gives \(E_{L'}\le B E_L\) after pullback and intersection with \(H^{d-1}\), so the same argument proves the last assertion. QED.

\textbf{Exact imported lattice source.} Martin Henk, \emph{Successive minima and lattice points}, arXiv:math/0204158v1, Theorem 1.3, printed p.2, and its proof in §3, printed pp.5--7, \url{https://arxiv.org/pdf/math/0204158}. The statement and complete proof were read; only the displayed upper-product inequality is imported here.

\textbf{What determinant is required.} On the real Hom space, let \(w_1,\ldots,w_r\) be a basis of normalized analytic vectors, with real Gram matrix \(G\). If the columns of \(B\) express an integral basis of \(\Lambda\) in these coordinates, then exactly \[
 C_\Lambda^2=\det(B^{\mathsf t}GB)=|\det B|^2\det G.
\tag{B10.3}
\] Replacing the analytic basis by an orthonormal one replaces \(B\) by \(G^{1/2}B\) and preserves this covolume. Rallis can contribute to \(G\) after the CM-period and measure comparisons are fixed; it does not identify \(B\). If \(\Gamma=\sum\mathbb Zw_j\) is commensurable with \(\Lambda\), then \[
 \frac{C_\Lambda}{C_\Gamma}
 =\frac{[\Gamma:\Gamma\cap\Lambda]}{[\Lambda:\Gamma\cap\Lambda]}.
\] Thus a proved integral analytic sublattice \(\Gamma\subset\Lambda\) with bounded Gram determinant already gives \(C_\Lambda\le C_\Gamma\); no saturation-index upper bound is then needed. If instead \(\Lambda\subset\Gamma\), its index contributes in the numerator and must be bounded. A uniform integer \(M\) with \(M\Gamma\subset\Lambda\) gives the sufficient bound \(C_\Lambda\le M^r\sqrt{\det G}\). This identifies the exact direction of the integral comparison that a theta-index theorem would have to supply.

\textbf{Variation and normalization.} The ranks \(r_i\) may grow when genuinely new cohomology appears at deeper source levels. Neither a rank-independent constant nor a uniform covolume is supplied by finite generation. For the fixed sublattice of old maps pulled through a degree-\(e\) cover of the source, pull back \(H\) as well. Then \(E_i\mapsto eE_i\), \(C_i\mapsto e^{r_i/2}C_i\), and a known \(\mu_i\) becomes \(e\mu_i\). Consequently \(R_i\mapsto\sqrt e R_i\), \(V\mapsto eV\), and \(\delta_F\mapsto e\delta_F\); the reduced-image bound (B10.1) is unchanged. The actual generic degree in that bound must be retained. Replacing it by the trivial lower bound one can lose an essential cancellation. This proposition gives a sufficient determinant criterion for producing a small witness from prescribed source lattices. It is not a necessary condition for intrinsic bounded witness degree.

\subsection{B11. A fixed-order CM Hecke example separates the Hom regulator from witness degree}\label{b11.-a-fixed-order-cm-hecke-example-separates-the-hom-regulator-from-witness-degree}

Let \[
 E=\mathbb C/(\mathbb Z+i\mathbb Z),\qquad
 E_m=\mathbb C/(\mathbb Z+mi\mathbb Z),\quad m\ge1,
\] with their degree-one polarizations. Then \[
 \operatorname{Hom}(E,E_m)=m\mathbb Z[i],\qquad
 \deg(m(a+bi))=m(a^2+b^2).
\tag{B11.1}
\] In particular the least positive degree grows with \(m\), despite fixed rank and fixed CM field.

\textbf{Proof.} A homomorphism is multiplication by a complex number \(\gamma\) with \(\gamma\mathbb Z[i]\subset\mathbb Z+mi\mathbb Z\). Write \(\gamma=x+mi y\) with \(x,y\in\mathbb Z\). The further requirement \(i\gamma\in\mathbb Z+mi\mathbb Z\) forces \(m\mid x\). Thus \(\gamma=m(a+bi)\), and the converse is immediate. The degree is the index of \(\gamma\mathbb Z[i]\) in the target lattice, namely \(|\gamma|^2/m=m(a^2+b^2)\). These elliptic curves are algebraic and have CM; multiplication by \(m\) gives a degree-\(m\) isogeny, so they lie in one polarized Hecke orbit. QED.

Here is a version with a fixed integral quadratic action and actual Weil witnesses. Put \(A_m=E_m^4\) with product principal polarization \(L_m\), and let \(\mathbb Z[i]\) act by sending \(i\) to \[
 J=\operatorname{diag}\left(
 \begin{pmatrix}0&-1\\1&0\end{pmatrix},
 \begin{pmatrix}0&-1\\1&0\end{pmatrix}\right).
\] This is an integral action for every \(m\), has Rosati involution \(J^\dagger=-J\), and has multiplicities \((2,2)\) on \(H^{1,0}\), so \(A_m\) is of Weil type for the same order. The diagonal isogeny \(E^4\to E_m^4\) given by multiplication by \(m\) commutes with this action and pulls the target polarization back to \(m\) times the source polarization. Thus the example stays in the corresponding polarized quadratic-action Hecke orbit.

Let \(Y=E^2\), \(H\) its product degree-one polarization, and consider \[
 f_m(x,y)=(mx,0,my,0),\qquad
 Z_m=E_m\times\{0\}\times E_m\times\{0\}\subset A_m.
\] The image has \(\deg_{L_m}Z_m=2\). Its Weil projection is nonzero: the extreme \(i\)-eigenline contains the nonzero \((2,2)\)-form \[
 (dz_1+i\,dz_2)\wedge(dz_3+i\,dz_4)
 \wedge(d\bar z_1+i\,d\bar z_2)
 \wedge(d\bar z_3+i\,d\bar z_4),
\] whose restriction to \(Z_m\) is \(dz_1\wedge dz_3\wedge d\bar z_1\wedge d\bar z_3\) and has nonzero integral. Pairing with this Weil form detects the nonzero projection, by the Poincaré self-adjoint projector in B6.

The full fixed-source Hom lattice is \[
 \operatorname{Hom}(E^2,A_m)=M_{4\times2}(m\mathbb Z[i]),
\] of rank \(16\). In the integral basis furnished by the real and imaginary matrix units, its energy metric is \(m I_{16}\): the polarization form on \(E_m\) is \(\frac{i}{2m}dz\wedge d\bar z\), so multiplication by \(m\gamma\) contributes \(m|\gamma|^2\) to the energy. Consequently its least nonzero energy is \(m\), its Gram determinant is \(m^{16}\), and its metric covolume is \(m^8\). Orthonormalizing the analytic coordinates transfers the factor \(\sqrt m\) to the integral coefficients and changes none of these intrinsic metric quantities.

For the selected two-dimensional map, however, \[
 f_m^*c_1(L_m)=mH,\quad
 E_{f_m}=2m,\quad V=H^2=2,\quad \delta_{f_m}=m^2,
\] and therefore \[
 \frac{E_{f_m}^2}{\delta_{f_m}V}=2=\deg_{L_m}Z_m.
\] All assertions follow by the lattice-index calculation, product intersection, and the displayed invariant form. This supplies a precise counterexample to any claim that finite generation or an orthonormal analytic basis uniformly bounds integral map energy on a CM Hecke orbit. At the same time it supplies bounded nonzero Weil witnesses. It therefore does \textbf{not} refute intrinsic bounded division; it shows that a raw Hom-regulator bound is sufficient but unnecessary, and that generic-degree cancellation is mathematically essential.

\subsection{B12. The exact geometric quantity is an infimum of normalized energies}\label{b12.-the-exact-geometric-quantity-is-an-infimum-of-normalized-energies}

For a polarized Weil-type abelian \(2n\)-fold \((A,L)\), let \(D_W(A,L)\) be the least degree of an effective integral \(n\)-cycle with nonzero Weil projection, or \(+\infty\) if none exists. Then \[
 D_W(A,L)=\inf_{(Y,f,H)}
 \frac{\bigl(f^*c_1(L)\cdot H^{n-1}\bigr)^n}
 {\delta_f (H^n)^{n-1}},
\tag{B12.1}
\] where \(Y\) runs over smooth projective \(n\)-folds, \(f:Y\to A\) over morphisms generically finite onto an image \(Z\) with nonzero Weil projection, \(\delta_f\) is the generic mapping degree, and \(H\) runs over integral ample classes on \(Y\). The infimum of the empty set is understood to be \(+\infty\).

\textbf{Proof.} The mixed-intersection inequality B10.2 and projection formula show that every expression on the right is at least \(\deg_L Z\ge D_W(A,L)\). If a witness exists, at least one irreducible component has nonzero Weil projection and no larger degree. Positive degrees are integers, so a least-degree irreducible witness \(Z\) exists. Take a projective resolution \(f:Y\to Z\subset A\), with \(\delta_f=1\), and put \(\beta=f^*c_1(L)\). This is integral, big and nef, and \(\beta^n=\deg_L Z=D_W(A,L)\). Fix an integral ample \(H_0\). The classes \(H_j=j\beta+H_0\) are integral ample. As \(j\to\infty\), \[
 \beta H_j^{n-1}=j^{n-1}D_W+O(j^{n-2}),\qquad
 H_j^n=j^nD_W+O(j^{n-1}).
\] For \(n=1\) the expression is already equality; for \(n\ge2\) the quotient in B12.1 tends to \(D_W\). This proves the reverse inequality and hence equality. QED.

This is an exact variational characterization, not an effective construction. It explains what a successful automorphic estimate would have to control: normalized energies for some witnesses, allowing their actual mapping degrees and source volumes. Requiring every theta coefficient or the entire Hom regulator to remain bounded is stronger than necessary. B10 provides one workable sufficient criterion with explicit constants. B11 shows why it is not an equivalence on a prescribed source. The infimum in B12 allows resolutions and polarizations of the unknown witnesses, so it does not close the open existence or uniformity problem by itself.

\subsection{B13. A fixed general-type surface bounds mapping degree to a simple target}\label{b13.-a-fixed-general-type-surface-bounds-mapping-degree-to-a-simple-target}

\textbf{Statement.} Let \(Y\) be a fixed smooth projective surface of general type, and let \(f:Y\to A\) have surface image \(Z\) in a simple abelian variety of dimension strictly greater than two. Then \[
 \delta_f\operatorname{vol}(K_{\widetilde Z})
 \le \operatorname{vol}(K_Y),\qquad
 \delta_f\le \operatorname{vol}(K_Y),
\tag{B13.1}
\] where \(\widetilde Z\) is any smooth projective model of \(Z\). In particular, if \(Y\) is a compact ball quotient, then \(\delta_f\le K_Y^2\). No smoothness assumption on \(Z\) and no unramifiedness assumption on \(f\) is required.

\textbf{Proof.} A proper positive-dimensional subvariety of a simple abelian variety is of general type by Ueno's theorem. Its Kodaira dimension is defined using a nonsingular model, so this applies even when \(Z\) is singular. The induced dominant rational map \(Y\dashrightarrow\widetilde Z\) has degree \(\delta_f\). The canonical-volume inequality for generically finite rational maps of varieties of general type gives the first inequality in (B13.1). One way to see the mechanism is to resolve the map, use the effective ramification divisor to compare the canonical class with the pullback of the target canonical class, and then use birational invariance, monotonicity, and finite-map scaling of volume. A smooth minimal model of a surface of general type has nef and big canonical divisor, whose square is a positive integer; its canonical volume is that square. Hence \(\operatorname{vol}(K_{\widetilde Z})\ge1\), proving the second inequality. For a compact ball quotient, \(K_Y\) is ample and its volume is \(K_Y^2\). QED.

\textbf{Exact imported sources.} Kenji Ueno, \emph{On Kodaira Dimensions of Certain Algebraic Varieties}, Proc. Japan Acad.47 (1971), pp.157--159, §3, Corollary after Proposition 2, printed p.158; the birational definition appears in §2, p.157. The note announces the theorem and says details will appear elsewhere; the whole note and the exact corollary were read, without claiming to have checked an absent full proof. Primary original: \url{https://www.jstage.jst.go.jp/article/pjab1945/47/2/47_2_157/_pdf}. The word ``proper'' is retained explicitly here, since the ambient abelian variety itself is not of general type. Lucio Guerra and Gian Pietro Pirola, \emph{On the finiteness theorem for rational maps on a variety of general type}, arXiv:0807.2980v2, §3.2, Proposition 3.2(2), printed pp.10--11, gives the canonical-volume inequality; its complete short proof was read. Primary: \url{https://arxiv.org/pdf/0807.2980}. Their later paper \emph{Parametrization of rational maps on a variety of general type, and the finiteness theorem}, arXiv:1203.2246v1, §1(b), printed pp.2--3, explicitly records the surface lower bound \(\epsilon_2=1\) and the degree-volume inequality. Primary: \url{https://arxiv.org/pdf/1203.2246}.

This is a strict refinement for a fixed ball-quotient source and a simple target: the unbounded generic-degree cancellation in B11 cannot occur there. It is not an energy estimate or a bound for all image degrees. For an étale source cover of degree \(e\), the canonical volume itself scales by \(e\), consistently with B3 and B10. The actual 032 target is generally a product, so simplicity is an additional hypothesis and must not be silently imposed. More generally the first inequality holds whenever the image has a general-type smooth model, without assuming the ambient target is simple.

\subsection{B14. Exact local bounds, and a fixed-level unit-input counterexample}\label{b14.-exact-local-bounds-and-a-fixed-level-unit-input-counterexample}

The local integrals in B1 cannot be bounded merely by fixing the unitary representation conductor, the target compact level, and the \(L^2\) norm of the Schwartz input. The following calculation makes the obstruction explicit and also gives a sharp bound with a stated support hypothesis.

\textbf{Statement at a split finite place.} Let \(F_v\) be nonarchimedean, with ring of integers \(\mathcal O_v\), uniformizer \(\varpi\), residue cardinality \(q\), additive Haar measure \(\operatorname{vol}(\mathcal O_v^3)=1\) and multiplicative Haar measure \(\operatorname{vol}(\mathcal O_v^\times)=1\). For an unramified compatible splitting, the rank-one member \(U(1)(F_v)=F_v^\times\) acts in \(\mathcal S(F_v^3)\) by \[
 \omega(a)=\rho(a)D(a),\qquad
 D(a)\phi(x)=|a|^{3/2}\phi(ax),
\] where \(\rho\) is a fixed unramified unitary character. Replacing \(a\) by \(a^{-1}\) in the model changes only the harmless phase convention below. Fix an unramified unitary source character \(\tau\). At the \(U(1)\to U(3)\) Rallis point, the diagonal local zeta integral, with a unit source vector, is \[
 Z_v(1;\phi)=\int_{F_v^\times}
       \langle\omega(a)\phi,\phi\rangle\tau(a)\,d^\times a.
\tag{B14.1}
\] For every integer \(M\ge1\) there is a \(GL_3(\mathcal O_v)\)-invariant Schwartz function \(\phi_M\) of \(L^2\) norm one such that \[
 Z_v(1;\phi_M)=M.
\tag{B14.2}
\] Thus \(Z_v^\#(1;\phi_M)\) is a fixed nonzero local Euler-factor ratio times \(M\).

\textbf{Proof.} The normalized linear model is Haan §2.3, equation (2.1), printed p.5; restricting the target \(GL_3\) action to its scalar center gives the displayed dilation, up to its fixed unitary character. The compatible center identification is Remark 2.1, p.4. Equation (3.3), printed p.8, expresses the Siegel--Weil section as the oscillator matrix coefficient times \(\Delta^{s-s_m}\). Here the source dimension is one, the target dimension is three and \(s_m=1\), so the extra factor is one. With the local source matrix coefficient this is exactly (B14.1). If a convention conjugates that coefficient, replace \(\tau\) by its inverse in the phase choice.

Put \(S_j=\varpi^j\mathcal O_v^3\setminus\varpi^{j+1}\mathcal O_v^3\) and \(e_j=\operatorname{vol}(S_j)^{-1/2}\mathbf1_{S_j}\) for \(j\in\mathbb Z\). The \(e_j\) are orthonormal, and \(D(\varpi)e_j=e_{j-1}\). Choose \(c=(\rho(\varpi)\tau(\varpi))^{-1}\), of modulus one, and set \[
 \phi_M=M^{-1/2}\sum_{j=0}^{M-1}c^j e_j.
\] It has norm one. Its support is contained in \(\mathcal O_v^3\) and it vanishes on \(\varpi^M\mathcal O_v^3\), so it is a Schwartz function. Every shell is \(GL_3(\mathcal O_v)\)-stable, and the splitting character is unramified; hence the full function is invariant under that same compact subgroup for every \(M\).

With the inner product linear in the first variable, \[
 \langle D(\varpi^k)\phi_M,\phi_M\rangle
 =c^k\left(1-\frac{|k|}{M}\right)
 \quad (|k|<M),
\] and the coefficient is zero for \(|k|\ge M\). The phases \(\rho(\varpi)^k\tau(\varpi)^kc^k\) cancel. Units contribute no phase and preserve every \(e_j\). Decomposing \(F_v^\times\) into \(\varpi^k\mathcal O_v^\times\), one obtains \[
 Z_v(1;\phi_M)=\sum_{|k|<M}\left(1-\frac{|k|}{M}\right)=M.
\] The sum is finite, so no continuation or interchange of divergent integrals is involved. The normalization in Haan Theorem 3.4 is multiplication by the fixed nonzero number \(L_{F_v}(3,\epsilon)/L_{E_v}(3/2,\operatorname{BC}(\tau)\gamma^3)\). At a split unramified place, writing \(\lambda\) for the first of the two inverse unitary Satake parameters, this number is \[
 \frac{|1-\lambda q^{-3/2}|^2}{1-q^{-3}}>0.
\] Thus the normalized local factors grow linearly too. QED.

\textbf{Sharp positive bound with radial support specified.} For \(\phi=\sum_{j\in J}a_j e_j\) supported on a finite set \(J\) of shells, absorb the same phases into \(a_j\). Summing its matrix coefficients gives \[
 Z_v(1;\phi)=\left|\sum_{j\in J}a_jc^{-j}\right|^2
 \le |J|\sum_{j\in J}|a_j|^2
 =|J|\|\phi\|_2^2.
\tag{B14.3}
\] This is Cauchy--Schwarz, and the \(M\)-shell example attains equality. Thus a support-size hypothesis gives an exact local bound, while the common group-theoretic level does not supply that support bound. The additive local constancy scale of \(\phi_M\) does become finer with \(M\); it has never been asserted fixed.

There is also a uniform bound using an explicit radial seminorm. For every \(\varepsilon>0\) and every integer center \(j_0\), weighted Cauchy--Schwarz gives \[
 Z_v(1;\phi)\le
 \bigl(2\zeta(1+\varepsilon)-1\bigr)
 \sum_{j\in\mathbb Z}(1+|j-j_0|)^{1+\varepsilon}|a_j|^2.
\tag{B14.5}
\] Indeed the sum of the inverse weights is \(1+2\sum_{k\ge1}(1+k)^{-1-\varepsilon}=2\zeta(1+\varepsilon)-1\). One may take the infimum over \(j_0\). Every radial Schwartz function has coefficients of this kind: only finitely many negative \(j\) occur, and local constancy near zero makes the positive tail decay as \(q^{-3j/2}\). Hence its coefficients are absolutely summable, the matrix-coefficient integral is absolutely convergent, and the exact radial formula and (B14.5) continue to hold for that tail. The estimate therefore converts the local-factor question into concrete support or weighted-seminorm control in this split radial model.

A pure dilation replaces the coefficient sequence by an index shift and a unit phase. It preserves both \(Z_v(1;\phi)\) and \(\|\phi\|_2\), and preserves the infimum of the seminorms over \(j_0\). Growth in (B14.2) comes from a coherent sum of many orthogonal shell vectors, rather than from translating one shell.

\textbf{Compact places.} If the local source torus \(T_v\) is compact, let its Haar volume be \(v_T\). The same diagonal section identity gives instead \[
 Z_v(1;\phi)=v_T\|P_{\tau^{-1}}\phi\|_2^2
 \le v_T\|\phi\|_2^2,
\tag{B14.4}
\] where \(P_{\tau^{-1}}=v_T^{-1}\int_{T_v}\tau(t)\omega(t)\,dt\) is the orthogonal character projector. This follows from compact-group averaging and unitarity. Multiplication by the explicitly retained local \(L\)-ratio gives the corresponding bound for \(Z_v^\#\). There is no positive lower bound for arbitrary inputs, since that projector may kill an input.

These are statements about the chosen theta inputs and local norm formulas. They do not imply that the intrinsic least degree of a geometric witness, or the intrinsic bounded-division complexity, grows. In particular the counterexample does not invalidate B8 or DIV1. It identifies precisely why an \(L\)-value estimate, even combined with unit input norms and fixed target level, is insufficient without additional control of the actual finite Schwartz data or an integral energy-lattice comparison.

\subsubsection{B14a. The split-place estimate for arbitrary angular inputs}\label{b14a.-the-split-place-estimate-for-arbitrary-angular-inputs}

The positive support and weighted bounds extend to every Schwartz input and to ramified unitary characters in the same split Schrödinger model. Let \(\mathcal H_0=L^2(S_0)\) and let \[
 (I_j b)(x)=q^{3j/2}b(\varpi^{-j}x)\mathbf1_{S_j}(x).
\] Then \(I_j:\mathcal H_0\to L^2(S_j)\) is unitary. Write \(\phi=\sum_j I_jb_j\), so \(\|\phi\|_2^2=\sum_j\|b_j\|_2^2\). Put \(\eta=\rho\tau\), \(u=\eta(\varpi)\), and let \(R(v)b(x)=b(vx)\) for \(v\in\mathcal O_v^\times\). The compact-unit operator \[
 P_\eta=\int_{\mathcal O_v^\times}\eta(v)R(v)\,d^\times v
\] is the orthogonal projector to the \(\eta^{-1}\) character. Then exactly \[
 Z_v(1;\phi)=
 \left\|\sum_{j\in\mathbb Z}u^jP_\eta b_j\right\|_{\mathcal H_0}^2.
\tag{B14a.1}
\] Consequently, for every \(\varepsilon>0\) and \(j_0\in\mathbb Z\), \[
 \begin{aligned}
 Z_v(1;\phi)
 &\le \bigl(2\zeta(1+\varepsilon)-1\bigr)
 \sum_j(1+|j-j_0|)^{1+\varepsilon}\|P_\eta b_j\|_2^2\\
 &\le \bigl(2\zeta(1+\varepsilon)-1\bigr)
 \int_{F_v^3}(1+|\nu(x)-j_0|)^{1+\varepsilon}|\phi(x)|^2\,dx,
 \end{aligned}
\tag{B14a.2}
\] where \(\nu(x)=\min_i v(x_i)\) for \(x\ne0\); the value assigned at zero is immaterial. If \(\phi\) is supported on at most \(M\) shells, then \(Z_v(1;\phi)\le M\|\phi\|_2^2\).

\textbf{Proof.} Normalized dilation satisfies \(D(\varpi^k)I_jb=I_{j-k}b\), while a unit acts by \(I_jR(v)b\). Integrating the unit part of (B14.1) replaces each inner product by \(\langle P_\eta b_j,b_{j-k}\rangle=\langle P_\eta b_j,P_\eta b_{j-k}\rangle\). Summing the valuation part and setting \(l=j-k\) gives \[
 \sum_{j,l}u^{j-l}\langle P_\eta b_j,P_\eta b_l\rangle
 =\left\|\sum_j u^jP_\eta b_j\right\|^2.
\] All interchanges are legitimate: compact support makes \(b_j=0\) for \(j\ll0\), and local constancy at zero gives \(b_j=q^{-3j/2}\phi(0)\mathbf1_{S_0}\) for \(j\gg0\). Hence \(\sum_j\|b_j\|<\infty\), and the double sum of absolute matrix-coefficient bounds is finite. Hilbert-space Cauchy--Schwarz with the weights from B14.5 proves the first estimate; the projection is a contraction and \(I_j\) is unitary, proving the integral expression and the second estimate. Ordinary finite-sum Cauchy--Schwarz proves the support bound. QED.

The unramified radial hypothesis was needed only to obtain the particularly simple fixed-\(GL_3(\mathcal O_v)\) counterexample in B14. It is not needed for this local bound. For a ramified character, the sufficiently deep constant tail may be annihilated by \(P_\eta\), which is already accounted for by the exact identity. Multiplying by the retained local Euler ratio gives the corresponding explicit estimate for the normalized zeta integral.

\subsection{B15. A linear bound in polarized isogeny degree on a Weil Hecke orbit}\label{b15.-a-linear-bound-in-polarized-isogeny-degree-on-a-weil-hecke-orbit}

\textbf{Setup.} Keep the connected polarized Weil-type family and fixed nonzero rational flat class \(t\) of B8. Fix one fibre \(A_0\) and effective integral cycles \(Z_0^+,Z_0^-\) such that \[
 N_0t_0=[Z_0^+]-[Z_0^-],\qquad
 \deg_{L_0}Z_0^\pm\le D_0,
\] with \(N_0\ge1\). Suppose a fibre \(A_x\) admits an \(\mathcal O\)-linear polarized isogeny \(\psi:A_x\to A_0\) with \[
 \psi^*c_1(L_0)=c\,c_1(L_x),\qquad c\in\mathbb Q_{>0},
\] and the same polarization type. The isogeny is not assumed to carry \(t_0\) to a rational multiple of \(t_x\).

\textbf{Statement.} Let \(b,\tau=bt,h,R,B_P,V_A,\delta_\Lambda\) be the fixed constants from B6--B8 and set \[
 \begin{aligned}
 C_N&=\frac{b h R^2}{(n!)^2 V_A\sqrt{\delta_\Lambda}},\\
 C_D&=\frac{2h\|\tau\|_q R B_P}{n!\sqrt{V_A}\sqrt{\delta_\Lambda}}.
 \end{aligned}
\] Then \(t_x\) has an integral signed-cycle representative \(N_xt_x=[Z_x^+]-[Z_x^-]\) satisfying \[
 \begin{aligned}
 1\le N_x&\le C_N D_0^2\deg\psi,\\
 \deg_{L_x} Z_x^\pm&\le C_D D_0^2\deg\psi.
 \end{aligned}
\tag{B15.1}
\] In particular its intrinsic bounded-division complexity is at most \(D_0^2\max(C_N,C_D)\deg\psi\). One may minimize this upper bound over all isogenies with the stated properties.

\textbf{Proof.} Polarization integration gives \(\deg\psi=c^{2n}\), and the codimension-\(n\) pullback calculation B4 gives \[
 \deg_{L_x}\psi^*Z_0^\pm
 =c^n\deg_{L_0}Z_0^\pm
 \le c^nD_0.
\] Flat pullback under an isogeny preserves effectiveness and integral coefficients. Because \(\psi\) is \(\mathcal O\)-linear, its pullback on \(H^1\) commutes with the two \(K\) eigenspaces, and hence induces an isomorphism \(W_0\to W_x\) on their extreme exterior powers. Equivalently it commutes with the fixed polynomial projector: \(e_x\psi^*=\psi^*e_0\). Therefore \[
 e_x[\psi^*Z_0^+]-e_x[\psi^*Z_0^-]
 =N_0\psi^*t_0\ne0.
\] At least one of these effective pullbacks is a nonzero Weil witness of degree \(D\le c^nD_0\). Apply B8 to that witness and to the prescribed target \(t_x\), using its two-dimensional extraction rather than assuming alignment of \(\psi^*t_0\) with \(t_x\). Substituting \(D^2\le c^{2n}D_0^2=D_0^2\deg\psi\) into the exact B8 constants gives (B15.1).

Uniformity follows from the family hypotheses, rather than from the isogeny: the order action makes \(W\) and its saturated integral lattice a local system, \(q\) is flat, \(\delta_\Lambda\) and \(\|bt\|_q\) are constant on the connected base, and \(V_A\) is fixed by polarization type. The fixed \(a\), \(P\), \(R\), \(h\) and \(B_P\) also do not vary. This proves the stated uniform constant. QED.

\textbf{Scalar special case.} If \(\psi^*t_0=r t_x\) for a rational scalar \(r\), then flatness gives \(\|t_0\|_q=\|t_x\|_q\), while Poincaré pairing under pullback gives \(\|\psi^*t_0\|_q^2=\deg\psi\,\|t_0\|_q^2\). Consequently \(r=\pm c^n\). Since \(c\) is rational and \(c^{2n}=\deg\psi\) is an integer, \(c\) is a positive integer. Pulling back the original signed representative, with the two signs interchanged when necessary, yields directly \[
 N_x=N_0c^n,\qquad
 \deg Z_x^\pm\le c^nD_0,
\] so the complexity is at most \(\max(N_0,D_0)\sqrt{\deg\psi}\). This is the scalar refinement of DIV4 in the fixed flat Weil-family normalization. The linear bound (B15.1) remains valid when the isogeny rotates the Weil plane and no such rational scalar exists.

This gives an unconditional quantitative upper bound once the fixed starting representative is supplied, for example at a fibre covered by an admitted Hodge-conjecture theorem. It does not assert a bound independent of isogeny degree and does not supply a lower bound on intrinsic complexity. Neither Dirichlet GRH nor Hecke GRH is used.

\subsection{B16. Exact Chow idempotence after the canonical degree projector}\label{b16.-exact-chow-idempotence-after-the-canonical-degree-projector}

The cohomological caution in B6 is essential for that elementary proof. For abelian varieties, an additional known motivic theorem does supply an exact Chow upgrade, with a fixed denominator. It is not obtained by inferring a Chow identity from Betti eigenvalues.

\textbf{Statement.} Work with rational Chow motives of an abelian variety \(A\) as in B6, or with relative Chow motives of an abelian scheme over a smooth connected quasi-projective complex algebraic base where the \(\mathcal O\) action is defined. Let \(\Pi_{2n}\) be the canonical Deninger--Murre degree projector. Then \[
 \mathcal E_W=\Pi_{2n}\circ\frac{P([a]^*)}{R}
\tag{B16.1}
\] is an idempotent correspondence modulo rational equivalence, whose Betti realization is the projection onto \(W\) in degree \(2n\) and zero in all other degrees. With \(M=[2]^*\), one may take \[
 \Pi_{2n}=\prod_{\substack{0\le j\le4n\\j\ne2n}}
              \frac{M-2^j}{2^{2n}-2^j}.
\tag{B16.2}
\] Thus a denominator for the displayed correspondence divides \[
 R D_n,\qquad
 D_n=\prod_{\substack{0\le j\le4n\\j\ne2n}}|2^{2n}-2^j|,
\] a fixed integer independent of the fibre and its extra CM endomorphisms.

\textbf{Imported motivic input.} Ben Moonen, \emph{On the Chow motive of an abelian scheme with non-trivial endomorphisms}, arXiv:1110.4264v2, §2.1 and Proposition 2.5--Corollary 2.6, PDF pp.7--8, states the exact canonical degree decomposition, its functoriality, the algebra homomorphism on \(h^1\), and the fact that \[
 r^{(i)}:\operatorname{End}^0(A)^{\mathrm{op}}
       \longrightarrow\operatorname{End}(h^i(A))
\] is a multiplicative homogeneous polynomial map of degree \(i\). The complete short proof of Corollary 2.6 was read. It uses the canonical exterior-power description, citing Künnemann, \emph{On the Chow motive of an abelian scheme} (1994), Theorem (3.3.1); Proposition 2.5 cites Kings (1998), Proposition 2.2.1. These motivic inputs are imported, not newly proved here. Primary source: \url{https://arxiv.org/pdf/1110.4264}.

\textbf{Proof.} The exact degree decomposition says that \(M\) acts by \(2^j\) on \(h^j(A)\) as a Chow-motive endomorphism. Consequently its interpolation polynomial in (B16.2) is exactly \(\Pi_{2n}\), not just its realization. Functoriality makes \(\Pi_{2n}\) commute with every \([a]^*\).

It remains to prove the polynomial relation for \(T=[a]^*\) on \(h^{2n}(A)\). Restrict the imported polynomial map \(r^{(2n)}\) to \(K\subset\operatorname{End}^0(A)\) and extend the coefficient field from \(\mathbb Q\) to \(K\). Under \(K\otimes_{\mathbb Q}K\simeq K\times K\), this becomes a multiplicative homogeneous polynomial map \[
 A(x,y)=\sum_{r=0}^{2n}E_r x^r y^{2n-r}
\] with coefficients \(E_r\) in the endomorphism algebra of \(h^{2n}(A)\) with \(K\) coefficients. Comparing coefficients in \(A(x,y)A(u,v)=A(xu,yv)\) gives \[
 E_rE_s=0\ (r\ne s),\qquad E_r^2=E_r.
\] The identity \(A(1,1)=1\) gives \(\sum_r E_r=1\). Evaluating at \((a,\bar a)\) yields \[
 T=\sum_{r=0}^{2n}a^r\bar a^{2n-r}E_r,
 \qquad \prod_{r=0}^{2n}(T-\lambda_r)=0
\] as exact identities of Chow-motive endomorphisms. Since scalar extension of a rational vector space is faithful, the rational polynomial relation holds already over \(\mathbb Q\).

The integral Bezout polynomial from B6 takes value \(R\) on the two extreme eigenvalues and zero on all others. Therefore \(P(T)/R=E_0+E_{2n}\) on \(h^{2n}(A)\) and is exactly idempotent there. Composing with \(\Pi_{2n}\) makes it zero on the other degrees, proving (B16.1) on the full Chow motive. Both factors are rational polynomials in algebraic pullback correspondences, and clearing their displayed denominators gives the asserted \(R D_n\) bound. Its realization agrees with B6 by the same extreme-eigenspace calculation. QED.

This does not assert that the cycle-class map is injective on the resulting Chow summand, or that every rational Weil class has an algebraic representative. It provides a canonical-degree Chow projector once the established abelian-motive decomposition is used. B8 and B15 already worked with the smaller cohomological polynomial \(P\) and do not require this additional motivic input.

\subsection{B17. Controlled local representatives and the remaining arithmetic comparison}\label{b17.-controlled-local-representatives-and-the-remaining-arithmetic-comparison}

This proposition refines B14--B14a. Arbitrarily many coherent radial shells can enlarge a local norm, but this presentation feature can be removed for a \textbf{specified existing} theta vector. The operation preserves that vector; it does not make its prescribed output small. A separate construction gives bounded local test vectors without prescribing their arithmetic realization.

\subsubsection{B17.1. An exact one-shell section}\label{b17.1.-an-exact-one-shell-section}

Use B14a's split model, measures and notation. In particular, \(\eta=\rho\tau\), \(u=\eta(\varpi)\), \(\mathcal H_0=L^2(S_0)\), and \[
 A_\eta\phi=\sum_j u^jP_\eta b_j,\qquad
 Z_\eta(\phi,\psi)=\langle A_\eta\phi,A_\eta\psi\rangle.
\] The second identity follows from the same absolutely convergent double sum as B14a, now with two inputs. Let \(\mathcal H_\eta^\infty\) be the locally constant functions on \(S_0\) transforming under scalar units by \(\eta^{-1}\). Define \[
 s_\eta(h)=I_0h,\qquad C_\eta=s_\eta A_\eta.
\] Then \(C_\eta\phi\) is a Schwartz function supported on the single shell \(S_0\), and \[
 \begin{aligned}
 A_\eta s_\eta&=1,& C_\eta^2&=C_\eta,\\
 A_\eta(\phi-C_\eta\phi)&=0,&
 Z_\eta(\phi-C_\eta\phi,\phi-C_\eta\phi)&=0,\\
 \|C_\eta\phi\|_2^2
 &=Z_\eta(C_\eta\phi,C_\eta\phi)
 =Z_\eta(\phi,\phi).
 \end{aligned}
\tag{B17.1}
\] For every \(\varepsilon>0\) its squared weighted shell norm from B14a, centered at zero, equals this same quantity. Moreover \(C_\eta\) commutes with \(GL_3(\mathcal O)\) in the linear target model, including a fixed smooth unitary determinant twist. It therefore preserves every compact subgroup level and compact type within this group.

\textbf{Proof.} Compact support gives \(b_j=0\) for \(j\ll0\). For some integer \(J\), local constancy at the origin gives \[
 b_j=q^{-3j/2}\phi(0)\mathbf1_{S_0}\quad(j\ge J).
\] Consequently \[
 A_\eta\phi
 =\sum_{j<J}u^jP_\eta b_j
 +\phi(0)P_\eta\mathbf1_{S_0}
       \frac{(u q^{-3/2})^J}{1-uq^{-3/2}}.
\tag{B17.2}
\] The first sum is finite, and the second is a locally constant function. Thus \(A_\eta\phi\in\mathcal H_\eta^\infty\). Extension by zero from the compact open shell is Schwartz. Its only nonzero shell coefficient is \(h\) at index zero, so \(A_\eta s_\eta h=P_\eta h=h\). This proves the section and idempotence identities. The bilinear form identity proves the radical and norm identities. For \(k\in GL_3(\mathcal O)\) the action \(V(k)\phi(x)=\kappa(\det k)\phi(k^t x)\) preserves each shell and acts identically on every shell coordinate. It commutes with scalar units and hence with \(P_\eta\), \(A_\eta\) and \(s_\eta\). The weighted norm assertion follows because only index zero remains. QED.

\textbf{Preservation of the actual global theta vector.} Assume the actual square-integrable rank-one theta lift and fixed splittings are in Haan's Theorem 3.4, printed p.8. For a pure tensor \(\Phi=\phi_v\otimes\Phi^v\), change one split finite factor to \(C_{\eta_v}\phi_v\). The difference has zero local raw zeta integral by (B17.1), and hence zero normalized integral: the normalizing Euler factors at \(3/2\) and \(3\) are finite and nonzero. Haan's \textbf{global diagonal norm identity} therefore gives \[
 \|\theta(\bar f,(\phi_v-C_{\eta_v}\phi_v)\otimes\Phi^v)\|_{\mathrm{Pet}}^2=0.
\] Positive definiteness and smoothness imply that this global automorphic theta vector is identically zero. Thus the replacement preserves the existing vector, not merely its norm. Apply this argument term by term for finite sums of pure tensors and sequentially at any finite set of places. It also applies componentwise to the finite-dimensional differential-form coefficient spaces in T6.11. In particular, if the original form realizes an integral algebraic homomorphism, its actual Betti periods and integral-Hom normalization are unchanged. This proof does not infer the kernel of a local theta quotient from a local norm alone.

Only finitely many places are changed; the distinguished unramified vectors are retained elsewhere. An infinite replacement by one-shell vectors need not be an adelic Schwartz function. If \(h=A_\eta\phi\ne0\), \(I_0h\) preserves the prescribed output and has norm \(\|h\|\); \(I_0(h/\|h\|)\) has unit output norm but changes the output by the scalar \(\|h\|^{-1}\). This distinction is essential for an arithmetic normalization.

\subsubsection{B17.2. Least-norm questions at split and compact source places}\label{b17.2.-least-norm-questions-at-split-and-compact-source-places}

For fixed nonzero \(h\in\mathcal H_\eta^\infty\), there is no least-\(L^2\)-norm Schwartz representative of its split output. Indeed \[
 \phi_{M,h}=\frac1M\sum_{j=0}^{M-1}u^{-j}I_jh
\] satisfies \[
 A_\eta\phi_{M,h}=h,\qquad
 \|\phi_{M,h}\|_2^2=\frac{\|h\|^2}{M}.
\tag{B17.3}
\] These are Schwartz functions with the same compact target level as \(h\). The identities follow from \(A_\eta I_jh=u^jh\) and orthogonality of different shells. Their squared norms tend to zero; zero cannot have output \(h\). Thus a minimal-input proposal requires an additional support, weighted, or arithmetic-lattice constraint.

At a compact source torus \(T\) with Haar volume \(v_T\), by contrast, \[
 P_{\tau^{-1}}=v_T^{-1}\int_T\tau(t)\omega(t)\,dt,
 \qquad Z_T(\phi,\phi)=v_T\|P_{\tau^{-1}}\phi\|^2.
\tag{B17.4}
\] The average is an orthogonal projector on smooth vectors. It commutes with the target action, and replacing \(\phi\) by \(P_{\tau^{-1}}\phi\) preserves the global theta vector by the preceding global norm argument. For a fixed projected output \(h\), any representative is \(h+(1-P)\phi\), an orthogonal sum, so \(h\) is the unique least-norm representative. This proves the assertion including its normalization.

\subsubsection{B17.3. Uniform analytic test vectors and the standard adelic tail}\label{b17.3.-uniform-analytic-test-vectors-and-the-standard-adelic-tail}

For any smooth unitary \(\eta\) the explicit function \[
 \phi_\eta(x)=\eta(x_1)^{-1}
       \mathbf1_{\mathcal O^\times\times\mathcal O\times\mathcal O}(x)
\tag{B17.5}
\] is Schwartz, has cyclotomic-integral values, and obeys \[
 A_\eta\phi_\eta=\phi_\eta|_{S_0},\qquad
 \|\phi_\eta\|_2^2=Z_\eta(\phi_\eta,\phi_\eta)=1-q^{-1}.
\tag{B17.6}
\] Thus its raw local zeta factor is between \(1/2\) and \(1\) independently of conductor. Normalization by \((1-q^{-1})^{-1/2}\) gives unit squared norm, unit raw zeta factor and unit squared weighted shell norm.

\textbf{Proof and level accounting.} The unit restriction of a smooth character has finite image, hence roots of unity. Scalar units act on (B17.5) by \(\eta^{-1}\); its support is one shell and has additive volume \(1-q^{-1}\). This proves (B17.6). For \(m\ge1\), \(k\in1+\varpi^mM_3(\mathcal O)\) preserves the chart and changes its first coordinate from \(x_1\) to \(x_1(1+\varpi^m z)\). Therefore \(m\ge\max(1,a(\eta))\) suffices for invariance in the untwisted model; with determinant twist \(\kappa\) it suffices that \(m\ge\max(1,a(\eta),a(\kappa))\). Conversely, for ramified \(\eta\) in the untwisted model, a nonzero \(A_\eta\)-image of a principal-level-\(m\) vector forces \(\eta\) to be trivial on \(1+\varpi^m\mathcal O\), by the scalar matrices in that subgroup. At level zero the same argument uses all units. Thus the conductor enters the level condition even though it does not enlarge this norm. These are Schwartz test vectors; no representation-theoretic newform identification is needed. QED.

For an unramified character the standard lattice vector \(\phi_0=\mathbf1_{\mathcal O^3}\) has \[
 Z_\eta(\phi_0,\phi_0)=\frac{1-q^{-3}}{|1-uq^{-3/2}|^2},\qquad
 \|\phi_0\|_{\mathrm{wt},2}^2
 =\frac{1+q^{-3}}{(1-q^{-3})^2}\le\frac{72}{49},
\tag{B17.7}
\] where the squared weight on shell \(j\) is \((1+|j|)^2\). Indeed \(b_j=q^{-3j/2}\mathbf1_{S_0}\) for \(j\ge0\), and zero otherwise. Summing the geometric series proves the first formula; differentiating it twice, or multiplying absolutely convergent series, gives \(\sum_{j\ge0}(j+1)^2t^j=(1+t)/(1-t)^3\), proving the second. Haan's standard unramified normalization gives normalized local factor one.

There is also a useful uniform restricted-product bound. Over a fixed number field \(F\), use the unit-normalized chart vector at any finite set of split places and the standard vector at all other split places. The product of their squared weighted norms is at most \[
 \prod_{v<\infty}\frac{1+q_v^{-3}}{(1-q_v^{-3})^2}
 =\frac{\zeta_F(3)^3}{\zeta_F(6)}.
\tag{B17.8}
\] Every factor on the right is greater than one. The Euler product converges absolutely, since its logarithm is \(O(q_v^{-3})\), and the equality follows from \(1+t=(1-t^2)/(1-t)\). Thus neither the number nor conductor of the selected split places forces growth in this purely analytic choice of local weighted norms. The statement prescribes no global integral-Hom class or Weil period.

\subsubsection{B17.4. Arithmetic coefficients retained by the section}\label{b17.4.-arithmetic-coefficients-retained-by-the-section}

The section preserving an existing vector need not preserve an integral lattice of Schwartz representatives. If \(U\subset\mathcal O^\times\) kills the unit action on a given vector and the character, normalized averaging is \[
 P_\eta b=\frac1{[\mathcal O^\times:U]}
       \sum_{v\in\mathcal O^\times/U}\eta(v)R(v)b.
\tag{B17.9}
\] This denominator occurs, for example, on \(b=\mathbf1_{U\times\mathcal O\times\mathcal O}\), whose projection is \([\mathcal O^\times:U]^{-1}\phi_\eta\). Shell coordinates also contain \(q^{-3j/2}\). Formula (B17.2) displays the exact geometric-tail denominator \(1-uq^{-3/2}\); the tail vanishes if \(\eta\) is ramified on units. For unramified \(\eta\) and \(u=1\), \[
 C_\eta(\mathbf1_{\mathcal O^3})
       =\frac{\mathbf1_{S_0}}{1-q^{-3/2}}.
\tag{B17.10}
\] When \(q=p^2\), the coefficient is the reduced rational number \(p^3/(p^3-1)\), with bounded real size and denominator \(p^3-1\). This is a local coefficient calculation, not a lower bound on geometric cycle complexity: the existing global vector is unchanged. The analogous compact-torus projector has an averaging denominator \([T:U]\) on any finite quotient supporting the vector.

\textbf{Main-gate consequence.} B17 removes arbitrary radial support complexity from an already specified rank-one theta vector without changing its level or integral realization. It also supplies analytic test inputs of uniformly bounded weighted norm. Neither assertion bounds the norm of the actual theta output realizing a primitive integral algebraic homomorphism. The remaining comparison is between that integral Hom lattice and the positive theta-output metric. B18 below removes a further small-period concern in the fixed generic rank-two Weil plane only after an integral middle-dimensional witness is realized. It does not supply the preceding lattice comparison.

The complete local derivation and a separate finite-ring model Fourier-period example are in \texttt{local\_newvector\_followup.md}, LN1--LN7. The model example is not identified with the geometric Weil period. Haan's imported exact formula is Theorem 3.4, with §2.3 equation (2.1), §3.1 equation (3.1), and §3.2 equation (3.3) fixing the local model and normalization. Source: \url{https://arxiv.org/pdf/1403.5061}.

\subsection{B18. A nonzero integral Weil witness has no arbitrarily small normalized Weil period}\label{b18.-a-nonzero-integral-weil-witness-has-no-arbitrarily-small-normalized-weil-period}

\textbf{Setup.} Keep the fixed generic rank-two Weil space \(W\subset H^{2n}(A,\mathbb Q)\), integral resultant projector \(Re=P([a]^*)\), saturated lattice \(\Lambda=W\cap H^{2n}(A,\mathbb Z)\) and positive integral form \[
 q(v,w)=(-1)^n\int_Av\smile w
\] from B6--B8. Choose \(\Omega\in W_\mathbb C\) in one of the two extreme \(K\)-eigenlines, normalized by \[
 q(\Omega,\Omega)=0,\qquad q(\Omega,\overline\Omega)=1.
\tag{B18.1}
\] The first identity is forced by the eigencharacters; the second can always be imposed by a positive real scaling. Either eigenline and any remaining phase give the same absolute values below. Let \(Z\) be an effective integral middle-dimensional cycle with \(w=e[Z]\ne0\).

\textbf{Statement.} There is the exact formula and uniform lower bound \[
 \left|\int_Z\Omega\right|
 =\frac{\|e[Z]\|_q}{\sqrt2}
 \ge\frac1R.
\tag{B18.2}
\] More precisely, \[
 R^2\left|\int_Z\Omega\right|^2
      =\frac12q(P[Z],P[Z])\in\mathbb Z_{>0}.
\tag{B18.3}
\] If \(D=\deg_L Z\) and \(V_A=\chi(L)\), B7 also gives \[
 \frac1R
 \le\left|\int_Z\Omega\right|
 \le\frac{D}{\sqrt2\,n!\sqrt{V_A}}.
\tag{B18.4}
\] For a proper generically finite map \(f:Y^n\to Z\subset A\) of generic degree \(\delta\), the corresponding source period satisfies \[
 \left|\int_Yf^*\Omega\right|
 =\delta\left|\int_Z\Omega\right|
 \ge\frac{\delta}{R}.
\tag{B18.5}
\] The constants are uniform in the connected family of B8 because \(R\), the generic Weil lattice and its integral pairing are fixed local-system data; no bound on the extra CM endomorphism field is required.

\textbf{Proof.} The extreme eigenlines are conjugate and paired nondegenerately by \(q\), while each is isotropic. They form all of \(W_\mathbb C\). Thus every real \(w\in W_\mathbb R\) has a unique expansion \(w=c\Omega+\bar c\overline\Omega\). Equation (B18.1) gives \[
 q(w,w)=2|c|^2,\qquad q(w,\Omega)=\bar c.
\] The projector \(e\) is Poincare-selfadjoint by B6 and fixes \(\Omega\). Consequently \[
 \int_Z\Omega
 =\int_A[Z]\smile\Omega
 =\int_Ae[Z]\smile\Omega
 =(-1)^nq(w,\Omega).
\] Taking absolute values proves the exact equality in (B18.2). The class \(v=Rw=P[Z]\) is a nonzero element of \(\Lambda\). Its square under \(q\) is a positive even integer. Positivity is Hodge--Riemann on \(W\). For evenness, the underlying real torus of \(A\) has integral cohomology ring \(\Lambda^*H^1(A,\mathbb Z)\), by the integral product computation for circles. Expand the positive even-degree class \(v\) in its ordered exterior-monomial basis. Each monomial has square zero, and two distinct monomials commute in this even degree, so every cross term in \(v^2\) occurs twice. Thus \(v^2\) is divisible by two in integral cohomology and \(q(v,v)\in2\mathbb Z_{>0}\). This proves (B18.3) and the lower bound. The B7 comass estimate gives the upper bound. Finally \(f_*[Y]=\delta[Z]\) as integral fundamental cycles, so the projection formula proves (B18.5). QED.

\textbf{Scope of the normalization.} This concerns the fixed generic two-dimensional Weil plane and a cycle already realized with integral coefficients. It does not concern a single full-CM embedding coefficient in the 032 construction, an arbitrary coordinate functional on a higher-dimensional CM space, or a complex-normalized theta form before arithmetic realization. Once such an integral witness exists, the scalar forcing this nonzero period to have absolute value one has absolute value at most \(R\) (at most \(R/\delta\) for (B18.5)). This eliminates a separate arbitrarily-small-period scaling factor at that stage. Complex period normalization need not preserve rational or integral cycle coefficients, and does not bound the original degree or integral-Hom energy. The positive bound in B8 still requires a degree bound for the effective witness.

\subsection{B19. An exact CM Hom covolume formula and a conductor bound for witnesses}\label{b19.-an-exact-cm-hom-covolume-formula-and-a-conductor-bound-for-witnesses}

This result controls an actual integral Hom lattice in a fixed simple CM isogeny class. It gives an arithmetic alternative to bounding the coefficients of a theta expansion. The relevant integer is the index of the largest lattice stable under the maximal CM order; it can vary along the Hecke orbit.

\subsubsection{B19.1. The exact metric covolume}\label{b19.1.-the-exact-metric-covolume}

Let \(E\) be a CM field of degree \(2g\), \(\Phi\) a fixed primitive CM type, \(\mathcal O=\mathcal O_E\), and \(D_E\) the absolute discriminant. Fix \[
 A_0=\mathbb C^g/\Phi(I_0),
\] where \(I_0\) is an invertible fractional \(\mathcal O\)-ideal, with ample polarization \(L_0\). Let \[
 A=\mathbb C^g/\Phi(\Lambda)
\] be an abelian variety in the same \(E\)-linear isogeny class, where \(\Lambda\subset E\) is a full integral lattice. Fix an ample polarization \(L\) on \(A\). Assume, as holds for polarizations of simple CM abelian varieties with full endomorphism field \(E\), that both Rosati involutions restrict to complex conjugation on \(E\). Put \[
 V_0=\chi(L_0),\qquad V=\chi(L),\qquad
 \Lambda^-=\{x\in\Lambda:\mathcal O x\subset\Lambda\},\qquad
 J=[\Lambda:\Lambda^-].
\tag{B19.1}
\] The lattice \(\Lambda^-\) is the largest \(\mathcal O\)-stable sublattice of \(\Lambda\). It has finite index because \(\Lambda\) is commensurable with \(\mathcal O\).

Let \(\omega_0,\omega\) be the invariant Kähler representatives of the two polarizations. For an \(E\)-linear homomorphism \(h:A_0\to A\), define the squared Hilbert--Schmidt norm \[
 Q_A(h)=\operatorname{tr}_{\omega_0}(h^*\omega).
\tag{B19.2}
\] The real rank of \(\operatorname{Hom}_E(A_0,A)\) is \(2g\). Its covolume in the positive metric \(Q_A\) is exactly \[
 \operatorname{covol}_{Q_A}\operatorname{Hom}_E(A_0,A)
 =2^{-g}\sqrt{|D_E|}\,\frac{V}{V_0}\,J.
\tag{B19.3}
\]

\textbf{Proof.} In the chosen CM coordinates, every \(E\)-linear rational homomorphism is multiplication by an element \(\alpha\in E\). It is an actual homomorphism precisely when \(\alpha I_0\subset\Lambda\). Thus its integral Hom lattice is the fractional \(\mathcal O\)-ideal \[
 I=(\Lambda:I_0)=\{\alpha\in E:\alpha I_0\subset\Lambda\}.
\] Invertibility of \(I_0\) gives \(I I_0=\Lambda^-\). Indeed \(I I_0\) is \(\mathcal O\)-stable and contained in \(\Lambda\), hence in \(\Lambda^-\). Conversely \(\Lambda^- I_0^{-1}\subset I\), and multiplying by \(I_0\) gives the reverse inclusion.

Use ordinary Euclidean measure \(\prod_{\sigma\in\Phi}dx_\sigma\,dy_\sigma\) on \(\mathbb C^g\). Write \(c(M)\) for the Euclidean covolume of a lattice \(M\subset E\). For fractional \(\mathcal O\)-ideals, multiplicativity of ideal norm gives \[
 c(I I_0)=\frac{c(I)c(I_0)}{c(\mathcal O)}.
\] This is also immediate by writing fractional ideals as products of prime ideals and using their lattice indices. Since \(c(\Lambda^-)=Jc(\Lambda)\), it follows that \[
 c(I)=J\,c(\Lambda)\frac{c(\mathcal O)}{c(I_0)}.
\tag{B19.4}
\] The standard ideal-covolume formula is Milne, \emph{Algebraic Number Theory}, version 3.08, Proposition 4.26, printed pp.79--80; the following determinant calculation verifies its convention here. The discriminant determinant in all \(2g\) embeddings has absolute value \(\sqrt{|D_E|}\). Passing each conjugate pair of embedding rows to its real and imaginary parts changes the absolute determinant by \(1/2\). Therefore \[
 c(\mathcal O)=2^{-g}\sqrt{|D_E|}.
\tag{B19.5}
\]

Rosati conjugation makes distinct CM eigenlines orthogonal. Hence \[
 \omega_0=\frac i2\sum_{\sigma\in\Phi}t_{0,\sigma}\,dz_\sigma\wedge d\bar z_\sigma,
 \qquad
 \omega=\frac i2\sum_{\sigma\in\Phi}t_\sigma\,dz_\sigma\wedge d\bar z_\sigma
\] for positive real coefficients. Explicitly, the adjoint relation for every \(a\in E\) forces an off-diagonal Hermitian coefficient between distinct embeddings to vanish. Integration over a fundamental domain gives \[
 V_0=c(I_0)\prod_\sigma t_{0,\sigma},\qquad
 V=c(\Lambda)\prod_\sigma t_\sigma.
\tag{B19.6}
\] Multiplication by \(\alpha\) has \[
 Q_A(\alpha)=\sum_\sigma\frac{t_\sigma}{t_{0,\sigma}}|\sigma(\alpha)|^2.
\] Changing from Euclidean metric to this real metric multiplies covolume by \(\prod_\sigma(t_\sigma/t_{0,\sigma})\): each complex coordinate contributes two real coordinates. Combining (B19.4)--(B19.6) proves (B19.3). In particular the period-lattice covolume cancels exactly. QED.

The analytic CM model used here follows from the rank-one \(E\)-module \(H_1(A,\mathbb Q)\) and the chosen CM type; an \(E\)-linear isogeny identifies the rational modules of all targets. Changing the identification by an element of \(E^\times\) scales \(\Lambda\) and \(\Lambda^-\) together, so \(J\) is unchanged. No comparison with an unspecified theta lattice has entered the calculation.

\subsubsection{B19.2. A small actual isogeny}\label{b19.2.-a-small-actual-isogeny}

Write \(\mu_A=\min_{0\ne h\in\operatorname{Hom}_E(A_0,A)}Q_A(h)\). The least nonzero squared Hom norm has the sharp order \[
 \begin{aligned}
 g\left(\frac{V}{V_0}J\right)^{1/g}
 &\le\mu_A\le K_E\left(\frac{V}{V_0}J\right)^{1/g},\\
 K_E&=\frac2\pi\bigl(g!\sqrt{|D_E|}\bigr)^{1/g}.
 \end{aligned}
\tag{B19.7}
\] A homomorphism \(h\) attaining the minimum is an actual isogeny and simultaneously satisfies \[
 J\le\deg h\le \frac{K_E^g}{g^g}\,J.
\tag{B19.8}
\] The lower bound \(\deg h\ge J\) holds for every nonzero integral \(E\)-linear homomorphism.

\textbf{Proof.} Any nonzero \(\alpha\in I\) acts invertibly on rational homology, so the resulting homomorphism is an isogeny, with \[
 \deg h=[\Lambda:\alpha I_0]\ge[\Lambda:\Lambda^-]=J,
\] because the \(\mathcal O\)-stable lattice \(\alpha I_0\) is contained in \(\Lambda^-\). Let \(\lambda_\sigma=(t_\sigma/t_{0,\sigma})|\sigma(\alpha)|^2\) be the positive eigenvalues of \(h^*\omega\) relative to \(\omega_0\). Then \[
 \sum_\sigma\lambda_\sigma=Q_A(h),\qquad
 \prod_\sigma\lambda_\sigma=\deg h\,\frac{V}{V_0}.
\] Arithmetic--geometric mean proves the lower bound in (B19.7).

Apply Minkowski's convex body theorem in real dimension \(r=2g\) to the closed ball in the metric \(Q_A\). The precise imported statement is J.S. Milne, \emph{Algebraic Number Theory}, version 3.08, Theorem 4.19, printed p.76: a compact symmetric convex set of volume at least \(2^r\) times the lattice covolume contains a nonzero lattice point. Since \(v_{2g}=\pi^g/g!\), (B19.3) gives \[
 \min_{h\ne0}Q_A(h)
 \le\left(\frac{2^{2g}\operatorname{covol}_{Q_A} I}{v_{2g}}\right)^{1/g}
 =K_E\left(\frac{V}{V_0}J\right)^{1/g}.
\] The minimum exists by positive definiteness and discreteness. For a minimizing \(h\), arithmetic--geometric mean in the other direction gives \[
 \deg h\le\frac{V_0}{V}\left(\frac{Q_A(h)}g\right)^g
 \le \frac{K_E^g}{g^g}J.
\] Together with the already proved lower bound, this proves all assertions. QED.

The cited statements and complete short proofs of Milne's Theorem 4.19 and Proposition 4.26 were read; these classical results are imported. Source: \url{https://www.jmilne.org/math/CourseNotes/ANT.pdf}.

\subsubsection{B19.3. Effective Weil witnesses and bounded division}\label{b19.3.-effective-weil-witnesses-and-bounded-division}

Now assume \(g=2n\), that the fixed generic imaginary quadratic field \(K\) of B6--B8 lies in \(E\), and that the CM type restricts to Weil signature \((n,n)\) on \(K\). Assume the order action and polarizations have the compatibility required there. Fix an effective integral cycle \(Z_0\subset A_0\) of dimension \(n\) with nonzero generic Weil projection and degree at most \(D_0\). Suppose the targets \(A\) belong to the fixed polarized Weil family of B8, with rational flat target class \(t\).

Then there is an effective integral nonzero Weil witness \(Z\) on each such target satisfying \[
 \deg_L Z
 \le D_0 K_E^n\left(\frac{V}{V_0}J\right)^{1/2}.
\tag{B19.9}
\] Using the explicit B8 constants \(C_N,C_D\) defined in B15, there is a signed integral representative \(N_A t_A=[Z_A^+]-[Z_A^-]\) with \[
 \begin{aligned}
 1\le N_A&\le C_N D_0^2 K_E^{2n}\frac{V}{V_0}J,\\
 \deg_L Z_A^\pm&\le C_D D_0^2 K_E^{2n}\frac{V}{V_0}J.
 \end{aligned}
\tag{B19.10}
\] These constants are independent of the endomorphism order of \(A\).

\textbf{Proof.} Choose the isogeny \(h\) from (B19.7) and take the effective integral pushforward cycle \(Z=h_*Z_0\), with its actual multiplicities. Positivity of Hermitian forms gives \[
 h^*\omega\le Q_A(h)\omega_0.
\] On each complex \(n\)-plane, the determinant is monotone for positive Hermitian forms; hence integration on the smooth loci of the components of \(Z_0\) gives \[
 \deg_L(h_*Z_0)
 =\int_{Z_0}(h^*\omega)^n
 \le Q_A(h)^n\deg_{L_0}Z_0.
\tag{B19.11}
\] Here no smoothness of the cycle or bound on the restriction's generic degree is needed. Since \(n/g=1/2\), (B19.7) proves (B19.9).

The isogeny is \(K\)-linear. Thus \(h^*\) commutes with the fixed \(K\) action, identifies the two generic Weil planes, and commutes with their polynomial projectors. Its pushforward is \(\deg h\,(h^*)^{-1}\) in cohomology, so it also commutes with those projectors and is an isomorphism on the Weil plane. Therefore \(e_A[h_*Z_0]=h_*e_0[Z_0]\ne0\). The effective cycle is a valid B8 witness. Substituting the square of (B19.9) into B8 gives (B19.10). The constants of B8 are flat-family constants as explained in B15; \(E\), \(V_0\) and \(D_0\) are fixed starting data. QED.

A fixed signed starting representative suffices: at least one of its two effective parts has nonzero Weil projection, and that part may be chosen for \(Z_0\). If the starting CM variety is first replaced by a fixed isogenous one with maximal-order action, transfer a fixed nonzero witness through that one fixed isogeny. Its resulting degree is then part of \(D_0\).

\subsubsection{B19.4. Conductor and polarized-isogeny comparisons}\label{b19.4.-conductor-and-polarized-isogeny-comparisons}

Let \[
 \mathcal R_A=\{a\in E:a\Lambda\subset\Lambda\},\qquad
 \mathfrak f_A=\{b\in\mathcal O:b\mathcal O\subset\mathcal R_A\}
\] be the CM endomorphism order and its conductor in \(\mathcal O\). Then \[
 J\le N_{E/\mathbb Q}(\mathfrak f_A).
\tag{B19.12}
\] Thus (B19.10) gives an unconditional upper bound linear in this conductor norm for the fixed CM isogeny class.

\textbf{Proof.} Put \(\Lambda^+=\mathcal O\Lambda\), a fractional \(\mathcal O\)-ideal. For \(b\in\mathfrak f_A\), \(a\in\mathcal O\) and \(x\in\Lambda\), the element \(ba\) belongs to \(\mathcal R_A\), so \(bax\in\Lambda\). Hence \(\mathfrak f_A\Lambda^+\subset\Lambda\), and this sublattice is \(\mathcal O\)-stable. It is therefore contained in \(\Lambda^-\). Consequently \[
 J\le[\Lambda:\mathfrak f_A\Lambda^+]
 \le[\Lambda^+:\mathfrak f_A\Lambda^+]
 =N_{E/\mathbb Q}(\mathfrak f_A),
\] as asserted. QED.

There is a useful comparison with B15 when both polarizations are principal. Suppose \(\psi:A\to A_0\) is an \(E\)-linear polarized isogeny with \(\psi^*\omega_0=c\omega\) for a positive rational \(c\). Then \[
 J\le\deg\psi.
\tag{B19.13}
\] Indeed identify \(\psi\) with multiplication by \(\alpha\) and put \(M=\alpha\Lambda\subset I_0\). Let \(M^\#\) denote the dual lattice for the integral alternating form of \(L_0\). Principal polarization gives \(I_0^\#=I_0\) and \(M^\#=c^{-1}M\). Write \(M^+=\mathcal O M\) and \(M^-\) for its largest \(\mathcal O\)-stable sublattice. The Rosati relation gives \[
 (M^+)^\#=(M^\#)^-=c^{-1}M^-.
\] Dualizing again and comparing finite indices yields \[
 J=[M:M^-]=[(M^-)^\#:M^\#]
   =[M^+:M]\le[I_0:M]=\deg\psi.
\] This proves (B19.13), and also identifies the core and hull indices exactly in this principal polarized situation. No corresponding comparison is asserted here without those polarization hypotheses.

\subsubsection{B19.5. Period-normalized analytic energy on a fixed source}\label{b19.5.-period-normalized-analytic-energy-on-a-fixed-source}

Suppose in addition that \(Y\) is a fixed smooth connected projective variety of dimension \(n\), with a fixed Kähler form \(\omega_Y\), and \(f_0:Y\to A_0\) is a morphism whose image has dimension \(n\). Let \(Z_0=f_{0*}[Y]\), with its actual generic multiplicity, and assume \(e_0[Z_0]\ne0\). Put \(g=2n\) as in B19.3, and choose \(h\) attaining the minimum in (B19.7). Define \[
 E_h=\int_Y(hf_0)^*\omega\,\omega_Y^{n-1}.
\] For each \(\sigma\in\Phi\) let \[
 a_\sigma=\int_Yf_0^*\left(\frac i2t_{0,\sigma}\,dz_\sigma\wedge d\bar z_\sigma\right)
                  \omega_Y^{n-1},\qquad
 a_{\min}=\min_\sigma a_\sigma,\quad a_{\max}=\max_\sigma a_\sigma.
\] These are fixed positive constants, and \[
 a_{\min}Q_A(h)\le E_h\le a_{\max}Q_A(h).
\tag{B19.14}
\] In particular the smallest energy of a nonzero integral map of this fixed form is bounded above and below by fixed positive multiples of \(J^{1/g}\).

Choose q-unit generic Weil forms \(\Omega_0,\Omega_A\) in corresponding extreme \(K\)-eigenlines as in B18, and put \[
 p_0=\left|\int_Yf_0^*\Omega_0\right|>0,\qquad
 p_h=\left|\int_Y(hf_0)^*\Omega_A\right|.
\] Then \[
 \begin{aligned}
 p_h&=\sqrt{\deg h}\,p_0\ge\sqrt J\,p_0,\\
 \frac{E_h}{p_h^{1/n}}
 &\le a_{\max}K_E\left(\frac{V}{V_0}\right)^{1/g}p_0^{-1/n}.
 \end{aligned}
\tag{B19.15}
\] Thus the period-normalized analytic energy is uniformly bounded in this fixed CM isogeny class and fixed-source construction.

\textbf{Proof.} Since \(A_0\) is simple and \(f_0\) is nonconstant, the induced Albanese homomorphism surjects onto \(A_0\). Pullback of invariant holomorphic one-forms is therefore injective. Every \(f_0^*dz_\sigma\) is nonzero, and its positive Kähler energy is strictly positive. This proves \(a_{\min}>0\). Diagonalization already used in B19.1 gives the exact identity \[
 E_h=\sum_\sigma a_\sigma\lambda_\sigma,
\] which proves (B19.14). Combining it with (B19.7) proves the asserted two-sided order for the minimum over all nonzero integral \(h\).

On middle cohomology, integration under the isogeny gives \[
 q_0(h^*u,h^*v)=\deg h\,q_A(u,v).
\] The \(K\)-linear isogeny preserves the chosen extreme eigenline, so \(h^*\Omega_A=\lambda\Omega_0\) for some \(\lambda\in\mathbb C^\times\). The q-unit normalization forces \(|\lambda|^2=\deg h\), proving the period identity. The lower bound uses (B19.8). Finally (B19.14) and (B19.7), together with \(g=2n\), give \[
 \frac{E_h}{p_h^{1/n}}
 \le\frac{a_{\max}K_E((V/V_0)J)^{1/g}}
          {p_0^{1/n}J^{1/(2n)}}
 =a_{\max}K_E(V/V_0)^{1/g}p_0^{-1/n}.
\] This proves (B19.15). QED.

For precision about the normalization, choose a polarization-unit invariant holomorphic coframe on \(A\). The sum of the fixed-source squared energies of its pullbacks is \(E_h\). Multiplying every pulled-back one-form by the positive scalar \(p_h^{-1/(2n)}\) makes the absolute value of the selected \(2n\)-fold generic Weil period equal to one and makes the summed squared energy equal to \(E_h/p_h^{1/n}\). The generic Weil form is a homogeneous wedge of \(2n\) one-forms and their conjugates; positive scalar scaling therefore has exactly this exponent.

When \(n=2\) and \(Y\) is a fixed compact unitary surface with the full rank-one theta spanning property of T6.9, this also gives uniformly controlled \textbf{analytic theta presentations}. Any one auxiliary finite cover required for that spanning statement is fixed at the outset, together with the pulled-back \(f_0\) and its multiplicity. Fix once and for all a finite basis \(\theta_1,\ldots,\theta_r\) of genuine theta one-forms on this source, with all its Schwartz data, source characters and levels. Define the positive energy Gram matrix \(G\) by \(\|\sum_k c_k\theta_k\|^2=c^*Gc\), and let \(\lambda_{\min}(G)>0\) be its least eigenvalue. Expansion of the normalized coordinate pullbacks in this basis gives \[
 \sum_{j,k}|c_{j,k}|^2
 \le \lambda_{\min}(G)^{-1}\frac{E_h}{p_h^{1/n}}.
\tag{B19.16}
\] Indeed the squared energy of each coefficient vector is \(c_j^*Gc_j\), and summing gives the numerator. A fixed linear section from this finite coefficient space to the finite collection of Schwartz presentations is bounded for every fixed Schwartz seminorm, including the weighted shell seminorms of B14a at the finitely many nonstandard places. Standard unramified inputs are kept outside that fixed set. Thus bounded period-normalized presentations exist in this explicitly fixed-source setting.

These normalized forms generally no longer arise from integral algebraic homomorphisms: the common real scaling changes their Betti lattice. The theorem supplies a positive answer to the analytic normalization question for this fixed reference construction, while retaining the arithmetic cost \(J\). It does not assert a uniform integral Hom norm, a uniform effective degree, a dense bounded-\(J\) locus, or a theorem for an arbitrary prescribed full-CM embedding coefficient in 032.

\textbf{Consequence and limit.} The Hom metric has an exact arithmetic size parameter; its determinant is not controlled solely by the finite Hecke \(L\)-value at \(3/2\). The integral Hom covolume is unbounded precisely when \(J\) is unbounded in this fixed-field, fixed-volume model, and the least nonzero squared Hom norm has the sharp order \(J^{1/g}\). Nonetheless a small actual isogeny and a Weil witness with degree \(O(\sqrt J)\) exist, giving complexity \(O(J)\). The elliptic lattices \(\mathbb Z+mi\mathbb Z\) of B11 already have \(J=m\) and minimum nonzero Hom squared norm \(m\), showing why a conductor-independent integral-Hom norm bound is unavailable in that example. Its image-cycle degree stays bounded, so this is still not a lower bound on intrinsic Weil complexity. No uniformity in \(E\) or its discriminant is asserted, and no GRH is used.

\subsection{Computation attached to B6}\label{computation-attached-to-b6}

\texttt{weil\_projector\_certificate.py} constructs the projector for \(K=\mathbb Q(\omega)\), \(\omega^2+\omega+1=0\), \(a=3+\omega\) (trace 5, norm 7), for \(n=2,3,4\), and checks exact polynomial idempotence modulo \(f_W f_R\). It also checks the norm-product resultant formula. This is a computation checking those examples, not a formal proof of B6 for all \(n\). For \(n=2\), for example, \[
 f_W=X^2-23X+2401,\quad f_R=X^3-126X^2+6174X-117649,
\] \[
 R=25\,729\,836\,300.
\] Reducing rational polynomial coefficients gives a smaller possible denominator \(476\,478\,450\) in this example. Minimality here means the least common denominator of this unique degree-\(\leq4\) interpolation polynomial, not a claim of a minimal denominator among all algebraic projector correspondences.

\subsection{Review of approach paths}\label{review-of-approach-paths}

\begin{enumerate}
\def\labelenumi{\arabic{enumi}.}
\tightlist
\item
  Tried direct GRH bound at the asserted edge. Refined by checking the actual Rallis dimensions: genuine torus-to-three-dimensional lift is at \(3/2\), and B1 is stronger and unconditional for the global factor.
\item
  Tried transfer from theta norm to degree. Direct identity is not justified; B5 supplies a valid inequality conditional on controlling all actual polarization coordinate pullbacks.
\item
  Tried conductor control along principal levels. B3 gives exact presentation growth and shows conductor is not level; no inference to intrinsic \(c(s)\) is made.
\item
  Tried use of bounded arbitrary projector norms. B9 refutes it. Replaced it by the fixed polynomial projector B6 in the generic \(K\)-action.
\item
  Tried treating the resulting rank-two coefficient index as a separate unbounded gate. B7 and B8 remove it once an effective degree bound is given, with explicit quadratic bounds. The residual gate is uniformly bounded effective witnesses with nonzero Weil component, or deformation of a full class that supplies them.
\end{enumerate}

\subsection{Suggested exact Grothendieck placement, to be independently source-checked by root}\label{suggested-exact-grothendieck-placement-to-be-independently-source-checked-by-root}

B6 concerns algebraic correspondences and primitive projectors in homological/numerical motives; B7 uses the Hodge standard positivity and the primitive Lefschetz decomposition; B8 concerns integral lattices underlying those motives together with proper cycle parameter spaces. The attached sketch points to cote 16 scans 44--50 for the standard-conjecture correspondence framework. No original archive scans are supplied here; this report does not claim to have read or corrected their transcription. Root must cite exact EGA/SGA/FGA and archive locators only after actual verification.

\subsection{Additional checked classical sources for B5 and B7}\label{additional-checked-classical-sources-for-b5-and-b7}

\begin{itemize}
\tightlist
\item
  Eduardo Cattani, \emph{Introduction to Kähler manifolds}, ICTP June 2010 preliminary notes, §5.3, Proposition 5.15 and Theorem 5.17, printed p.38: the primitive Hodge-star formula and Hodge--Riemann positivity. Example 5.19 and equation (5.23), printed pp.39--40, specialize to the surface Hodge inequality. The displayed statements and the adjoining positivity proof were read. Source: \url{https://people.maths.ox.ac.uk/joyce/Nairobi2019/CattaniKahlerManifolds.pdf}. These are classical imported results; this report does not claim an independent proof of the full Hodge--Riemann theorem.
\item
  Jean-Pierre Demailly, \emph{Complex Analytic and Differential Geometry}, version 21 June 2012, Chapter III, §1.D, equation (1.24) and Wirtinger's inequality (1.25), printed pp.136--137: analytic-cycle volume and its linear-algebra proof. Source: \url{https://people.math.harvard.edu/~demarco/Math274/Demailly_ComplexAnalyticDiffGeom.pdf}. Demailly uses \(\omega_D=i\sum h_{jk}dz_j\wedge d\bar z_k\) and the real metric \(\operatorname{Re}h\), giving \(\omega_D^p/(2^p p!)\). For the same real metric our convention is \(\omega=\omega_D/2\), so this is precisely \(\omega^p/p!\) as used in B7. The formula and proof were read and the factor of two checked.
\end{itemize}
